%% =========================================================
%% Chapter: Marginal surprisal moments under a global gap
%% Source: OpenAI, "A two-dimensional area law from a global spectral gap"
%%         (September 24, 2026), Lemma 3.1, 02-initial.tex lines 34-207
%% =========================================================

\chapter{Marginal Surprisal Moments Under a Global Gap}
\label{ch:marginal_tails}

For a gapped ground vector of a sum of bounded terms on a bipartite space,
the negative logarithm of a marginal eigenvalue, sampled with that
eigenvalue as its probability, concentrates around the entropy of the
marginal. This chapter proves the moment and tail bounds
of~\cite[Lemma~3.1]{OpenAI2026AreaLaw} for a bipartition $\mathcal
    H_B\otimes\mathcal H_{B^c}$. Every term acts on one factor only, except the
crossing terms, each of which is a sum of at most $d_i^2$ products
$c_\alpha\otimes d_\alpha$ with $\lVert c_\alpha\rVert\lVert d_\alpha\rVert\le c_0$.
A term supported on sites of total dimension $d_i$ has such a
decomposition through the matrix units on its portion in $B$.

Throughout, $p$ is a probability vector, $K$ takes the value $-\log p_j$
with probability $p_j$, $M(u)=\mathbb E\,\mathrm e^{uK}=\sum_{p_j>0}p_j^{1-u}$, and
$S=\mathbb E K$ is the Shannon entropy.

\section{Scalar estimates}
\label{sec:marginal_tails_scalar}

\begin{lemma}[Quadratic real part from three lines]
    \label{lem:strip_quadratic}
    \lean{Complex.norm_le_exp_one_mul_of_strip,
        Complex.abs_re_sub_le_of_norm_le_closedBall,
        Complex.abs_re_sub_le_of_strip}
    \leanok
    Let $f$ be entire with $f(-x)=\overline{f(x)}$ for real $x$,
    $|f(\mathrm iy)|\le c$ for real $y$, and $|f(z)|\le cm$ for
    $|\operatorname{Re}z|\le1/4$, where $m\ge1$. If $\ell\ge1$ and
    $\log m\le2\ell$, then $|f(z)|\le\mathrm ec$ for
    $|\operatorname{Re}z|\le1/(8\ell)$, and
    $|\operatorname{Re}(f(x)-f(0))|\le128\mathrm ec\ell^2x^2$ for real
    $|x|\le1/(8\ell)$.
\end{lemma}
\begin{proof}\leanok
    The three-lines theorem on the strips between real parts $0$ and $\pm1/4$
    gives $|f(z)|\le c\,m^{4|\operatorname{Re}z|}\le\mathrm ec$. Put
    $k(z)=f(z)+f(-z)-2f(0)$, so that $k(0)=k'(0)=0$ and $|k|\le4\mathrm ec$
    on the disk of radius $\delta=1/(8\ell)$. The function $k(z)/z^2$ is
    entire, and the maximum principle bounds it by $4\mathrm ec/\delta^2$ on
    the disk. Since $\operatorname{Re}f(-x)=\operatorname{Re}f(x)$, the
    quantity $\operatorname{Re}(f(x)-f(0))$ equals $\operatorname{Re}k(x)/2$.
    Only the real part of the linear Taylor coefficient vanishes.
\end{proof}

\begin{lemma}[Halving recurrence and exponential tails]
    \label{lem:surprisal_halving}
    \lean{Entropy.hasDerivAt_log_surprisalMoment, Entropy.le_mul_add_of_halving,
        Entropy.log_surprisalMoment_le_of_gap, Entropy.surprisalTail_le}
    \leanok
    The function $L=\log M$ satisfies $L(0)=0$ and $L'(0)=S$. If $K_0\ge0$,
    $K_0u^2\le1/2$ and $(1-K_0u^2)M(u)\le M(u/2)^2$ for all $|u|\le r$, then
    $L(u)\le uS+4K_0u^2$ for $|u|\le r$. If $L(\pm r)\le\pm rS+Cr^2$ with
    $r>0$, then $\Pr\{|K-S|>w\}\le2\mathrm e^{Cr^2}\mathrm e^{-rw}$.
\end{lemma}
\begin{proof}\leanok
    The recurrence gives $L(u)\le2L(u/2)+2K_0u^2$ because
    $-\log(1-v)\le2v$ for $0\le v\le1/2$. Iterating, $2^jL(u/2^j)\to uS$, and
    the geometric series of quadratic terms sums to at most
    $4K_0u^2$. The tail is exponential Markov at $u=\pm r$.
\end{proof}

\section{Modular conjugation}
\label{sec:marginal_tails_modular}

Let $\Theta$ be a unit vector whose marginal on $\mathcal H_B$ is
$\operatorname{diag}(s)$, with rows $\Theta_j$. For an operator $X$ put
$f(z)=\sum_{j,k}\mathrm e^{z(\log s_j-\log s_k)}\langle\Theta_j,X_{jk}\Theta_k\rangle$,
which is $\langle\Theta,\sigma^zX\sigma^{-z}\Theta\rangle$ on the support of
$\sigma=\operatorname{diag}(s)$.

\begin{lemma}[Local conjugation estimate]
    \label{lem:modular_conjugation_quadratic}
    \lean{Entropy.modularExpectation_zero, Entropy.norm_modularExpectation_mul_I_le,
        Entropy.modularExpectation_neg, Entropy.norm_modularExpectation_kronecker_le,
        Entropy.abs_re_modularExpectation_sub_le}
    \leanok
    \uses{lem:strip_quadratic}
    One has $f(0)=\langle\Theta,X\Theta\rangle$ and $|f(\mathrm iy)|\le\lVert X\rVert $,
    and $f(-x)=\overline{f(x)}$ for Hermitian $X$. For $X=c\otimes d$ and
    $|\operatorname{Re}z|\le1/2$, $|f(z)|\le2\lVert c\rVert \lVert d\rVert $. Consequently,
    if $X$ is Hermitian, $\lVert X\rVert \le c_0$ and $X$ is a sum of $N$ products
    with $\lVert c_\alpha\rVert\lVert d_\alpha\rVert\le c_0$, then
    $|\operatorname{Re}(f(x)-f(0))|\le128\mathrm ec_0\ell^2x^2$ whenever
    $\ell\ge1$, $\log(2N)\le2\ell$ and $|x|\le1/(8\ell)$.
\end{lemma}
\begin{proof}\leanok
    On the imaginary axis $f$ is an expectation in a vector with the same
    norm. For a product, write $\Theta_j=\sqrt{s_j}\,\omega_j$ with orthonormal
    $\omega_j$. The weights satisfy
    $|s_j^{1/2+z}s_k^{1/2-z}|\le s_j+s_k$, and Cauchy--Schwarz with Bessel's
    inequality bounds every row and column sum of
    $|c_{jk}\langle\omega_j,d\,\omega_k\rangle|$ by $\lVert c\rVert \lVert d\rVert $. The last
    claim is Lemma~\ref{lem:strip_quadratic} with $m=2N$.
\end{proof}

\section{The moment and tail bounds}
\label{sec:marginal_tails_main}

\begin{lemma}[Marginal tilt]
    \label{lem:marginal_tilt}
    \lean{Entropy.norm_tilt, Entropy.inner_tilt,
        Entropy.modularExpectation_tilt_tiltExponent,
        Entropy.modularExpectation_kronecker_one, Entropy.modularExpectation_one_kronecker,
        Entropy.re_inner_tilt_sub_le, Entropy.gap_surprisalMoment}
    \leanok
    \uses{lem:modular_conjugation_quadratic}
    Suppose the marginal of the ground vector $\Psi$ is $\operatorname{diag}(p)$.
    The vector $\Theta=F\Psi/\sqrt{M(u)}$, $F=\operatorname{diag}(p^{-u/2})$ on
    the support, is a unit vector with marginal proportional to $p^{1-u}$ and
    $\langle\Psi,\Theta\rangle=M(u/2)/\sqrt{M(u)}$. At
    $z_u=-u/(2(1-u))$ its modular function is $\langle F^2\Psi,X\Psi\rangle/M(u)$,
    and one-sided terms have constant modular function. For
    $|u|\le1/2$ and $|u|\le1/(8\log(\mathrm ed_i))$ at every crossing term,
    $\langle\Theta,(\mathsf H-E_0)\Theta\rangle\le128\mathrm ec_0\mathcal Bu^2$,
    and the gap gives $(1-128\mathrm e\vartheta\mathcal Bu^2)M(u)\le M(u/2)^2$.
\end{lemma}
\begin{proof}\leanok
    Summing the value at $z_u$ over all terms gives
    $\langle F^2\Psi,\mathsf H\Psi\rangle/M(u)=E_0$. The excitation energy is
    therefore the sum of $\operatorname{Re}(f_i(0)-f_i(z_u))$ over the crossing
    terms, each bounded by Lemma~\ref{lem:modular_conjugation_quadratic}
    since $|z_u|\le|u|$. The gap inequality
    $\mathsf H-E_0\ge g_0(1-\ket\Psi\bra\Psi)$ then bounds
    $1-|\langle\Psi,\Theta\rangle|^2$.
\end{proof}

\begin{theorem}[Marginal moment and tail bounds, bipartite form]
    \label{thm:marginal_tails}
    \lean{Entropy.log_surprisalMoment_eigenvalues_le,
        Entropy.surprisalTail_eigenvalues_le,
        Entropy.cutBasisChange_conj_sum_kronecker}
    \leanok
    \uses{lem:marginal_tilt, lem:surprisal_halving}
    Let $\mathsf H=\sum_ih_i$ on $\mathcal H_B\otimes\mathcal H_{B^c}$ with
    Hermitian terms, $\lVert h_i\rVert\le c_0$, $c_0\ge0$, each term acting on one
    factor except the crossing terms $i\in\mathcal I_\times$, each a sum of
    at most $d_i^2$ products with factor norms at most $c_0$, $d_i\ge1$. Let
    $\Psi$ be a unit vector with $\mathsf H\Psi=E_0\Psi$ and
    $\mathsf H-E_0\ge g_0(1-\ket\Psi\bra\Psi)$, $g_0>0$. Let $\rho$ be its
    marginal on $\mathcal H_B$, $p$ the eigenvalues of $\rho$, $S=S(\rho)$,
    $\mathcal B=1+\sum_{i\in\mathcal I_\times}\log^2(\mathrm ed_i)$ and
    $\vartheta=c_0/g_0$. Then
    $\log\mathbb E\,\mathrm e^{uK}\le uS+512\mathrm e\vartheta\mathcal Bu^2$ for
    $|u|\le1/(32\sqrt{(1+\vartheta)\mathcal B})$, and for every real $w$
    \[
        \Pr\{|K-S|>w\}\le\min\Bigl\{1,\,2\mathrm e^{\mathrm e/2}
        \exp\Bigl(-\frac{w}{32\sqrt{(1+\vartheta)\mathcal B}}\Bigr)\Bigr\}.
    \]
\end{theorem}
\begin{proof}\leanok
    Conjugating by $U^*\otimes1$, where $U$ diagonalizes $\rho$, preserves
    one-sided terms, product decompositions with their factor norms, the
    eigenvalue equation and the gap inequality, and makes the marginal
    diagonal. On the admissible interval every halving stays admissible,
    $|u|\le1/32$, $\log(\mathrm ed_i)\le\sqrt{\mathcal B}$ and
    $128\mathrm e\vartheta\mathcal Bu^2\le1/2$, so
    Lemmas~\ref{lem:marginal_tilt} and~\ref{lem:surprisal_halving} give the
    moment bound. At $u=\pm1/(32\sqrt{(1+\vartheta)\mathcal B})$ the
    quadratic term is at most $\mathrm e/2$, which gives the tail.
\end{proof}

\section{Designated supports}
\label{sec:marginal_tails_supports}

Let $\mathcal H=\bigotimes_{v\in V}\C^{n_v}$. An operator is supported on
$D\subseteq V$ when it acts as the identity outside $D$: its matrix
elements vanish between configurations that differ outside $D$ and do not
depend on the common values there.

\begin{lemma}[Matrix-unit decomposition across a cut]
    \label{lem:supported_product_decomposition}
    \lean{Entropy.IsSupportedOn.isOneSided_of_subset,
        Entropy.IsSupportedOn.isOneSided_of_disjoint,
        Entropy.IsSupportedOn.hasProductDecomposition}
    \leanok
    Let $h$ be supported on $D$ and let $B\sqcup B^c=V$. If $D\subseteq B$ or
    $D\subseteq B^c$, then $h$ acts on one factor of
    $\mathcal H_B\otimes\mathcal H_{B^c}$. In general $h$ is a sum of at most
    $d^2$ products $c_\alpha\otimes d_\alpha$ with $\lVert c_\alpha\rVert\le1$
    and $\lVert d_\alpha\rVert\le\lVert h\rVert$, where
    $d=\prod_{v\in D}n_v$.
\end{lemma}
\begin{proof}\leanok
    The factor $c_\alpha$ is a matrix unit on $D\cap B$ tensored with the
    identity on $B\setminus D$, written as $S_aS_b^*$ for isometries
    $S_a$. The factor $d_\alpha$ is the compression $T_a^*hT_b$ along the
    isometries $T_a\colon\psi\mapsto e_a\otimes\psi$. The identity
    $h=\sum c_\alpha\otimes d_\alpha$ follows entrywise from the support
    property, and there are $(\prod_{v\in D\cap B}n_v)^2\le d^2$ terms.
\end{proof}

\begin{theorem}[Marginal moment and tail bounds]
    \label{thm:marginal_tails_supports}
    \lean{Entropy.log_surprisalMoment_cut_le, Entropy.surprisalTail_cut_le}
    \leanok
    \uses{thm:marginal_tails, lem:supported_product_decomposition}
    Let $\mathsf H=\sum_{i\in\mathcal I}h_i$ on $\mathcal H$ with $\mathcal I$
    finite, each $h_i$ Hermitian with $\lVert h_i\rVert\le c_0$, $c_0\ge0$, and
    supported on a designated set $D_i$. Let $\Psi$ be a unit vector with
    $\mathsf H\Psi=E_0\Psi$ and $\mathsf H-E_0\ge g_0(1-\ket\Psi\bra\Psi)$,
    $g_0>0$. Fix $B\sqcup B^c=V$, let $\mathcal I_\times$ be the indices whose
    support meets both sides, $\rho=\Tr_{B^c}\ket\Psi\bra\Psi$,
    $d_i=\prod_{v\in D_i}n_v$,
    $\mathcal B=1+\sum_{i\in\mathcal I_\times}\log^2(\mathrm ed_i)$ and
    $\vartheta=c_0/g_0$. Then the conclusions of
    Theorem~\ref{thm:marginal_tails} hold:
    $\log\mathbb E\,\mathrm e^{uK}\le uS+512\mathrm e\vartheta\mathcal Bu^2$
    for $|u|\le1/(32\sqrt{(1+\vartheta)\mathcal B})$, and the stated tail bound
    holds for every real $w$.
\end{theorem}
\begin{proof}\leanok
    Pass to cut coordinates. The operator norm, Hermiticity, the eigenvalue
    equation and the gap inequality are preserved by the relabelling.
    Lemma~\ref{lem:supported_product_decomposition} supplies the one-sided
    terms and the product decompositions of the crossing terms. Every local
    dimension is positive, since $\Psi\ne0$. Apply
    Theorem~\ref{thm:marginal_tails}.
\end{proof}

\begin{lemma}[Typical sets]
    \label{lem:typical_set}
    \lean{Entropy.one_sub_surprisalTail_le_typicalMass, Entropy.exp_le_of_mem_typicalSet}
    \leanok
    Let $E$ be the set of indices with $p_i>0$ and $|-\log p_i-S|\le w$, and
    let $z=\sum_{i\in E}p_i$. Then $z\ge1-\Pr\{|K-S|>w\}$, and
    $\mathrm e^{-S-w}\le p_i\le\mathrm e^{-S+w}$ for $i\in E$.
\end{lemma}
\begin{proof}\leanok
    The indices outside $E$ with positive weight are exactly those in the
    tail event. The weight bounds are the definition of $E$.
\end{proof}

\section{Normalized filters}
\label{sec:marginal_tails_filters}

\begin{lemma}[Trace Hölder inequality]
    \label{lem:trace_holder}
    \lean{Matrix.trace_mul_eq_sum_mul_normSq, Matrix.re_trace_mul_le_of_spectral,
        Matrix.PosSemidef.re_trace_mul_le}
    \leanok
    For positive semidefinite $A,B$ with eigenvalues $a_i,b_j$ and
    $1/p+1/q=1$, $p,q>1$,
    $\operatorname{Re}\Tr(AB)\le(\sum_ia_i^p)^{1/p}(\sum_jb_j^q)^{1/q}$.
\end{lemma}
\begin{proof}\leanok
    With $A=U\operatorname{diag}(a)U^*$ and $B=V\operatorname{diag}(b)V^*$,
    $\Tr(AB)=\sum_{ij}a_ib_j|W_{ij}|^2$ for the unitary $W=U^*V$, whose squared
    moduli are doubly stochastic. Apply the scalar Hölder inequality to
    $a_i|W_{ij}|^{2/p}$ and $b_j|W_{ij}|^{2/q}$.
\end{proof}

\begin{lemma}[Norm of a normalized filter]
    \label{lem:filter_norm}
    \lean{Entropy.norm_kronecker_one_apply_sq,
        Entropy.norm_kronecker_one_apply_sq_le_rpow,
        Entropy.norm_kronecker_one_apply_sq_le_exp,
        Entropy.sum_eigenvalues_partialTraceRight_eq_one,
        Entropy.rpow_one_sub_le_exp_of_log_le}
    \leanok
    \uses{lem:trace_holder, thm:marginal_tails}
    Let $L\ge0$ act on $\mathcal H_B$ with $\Tr L^{2/a}=1$, $0<a<1$. For a
    unit vector with marginal $\rho$, $\lVert(L\otimes1)\Psi\rVert^2=\Tr\rho
        L^2\le(\Tr\rho^{1/(1-a)})^{1-a}=M(-a/(1-a))^{1-a}$. Under the hypotheses
    of Theorem~\ref{thm:marginal_tails}, if $a\le1/2$ and
    $a/(1-a)\le1/(32\sqrt{(1+\vartheta)\mathcal B})$, then
    $\lVert(L\otimes1)\Psi\rVert^2\le\exp(-aS(\rho)+1024\mathrm e\vartheta
        \mathcal Ba^2)$.
\end{lemma}
\begin{proof}\leanok
    Apply Lemma~\ref{lem:trace_holder} with exponents $1/(1-a)$ and $1/a$ to
    $\rho$ and $L^2$, then the moment bound at $u=-a/(1-a)$, and use
    $1/(1-a)\le2$.
\end{proof}

\begin{lemma}[Clipped eigenvalues and conjugation cost]
    \label{lem:clipped_filter_cost}
    \lean{Entropy.abs_log_max_sub_log_max_le, Entropy.abs_log_clipped_sub_le,
        Entropy.cosh_sub_one_le, Entropy.sqrt_mul_cosh_sub_one_le}
    \leanok
    For $\varepsilon,\lambda>0$, $a\ge0$ and positive $p_i,p_k$, the numbers
    $l_i=\max\{\varepsilon,p_i/\lambda\}^{a/2}$ satisfy
    $|\log(l_i/l_k)|\le\frac a2|\log(p_i/p_k)|$. If $a\le1/2$, $p,q>0$ and
    $|t|\le\frac a2|\log(p/q)|$, then
    $\sqrt{pq}\,(\cosh t-1)\le2a^2(p+q)$.
\end{lemma}
\begin{proof}\leanok
    A maximum with a common floor contracts differences of logarithms. For
    the cost, write $p=\mathrm e^{m+z}$, $q=\mathrm e^{m-z}$, so that
    $p+q=2\sqrt{pq}\cosh z$ and $|t|\le a|z|\le|z|/2$. Then
    $\cosh t-1\le\frac{t^2}2\cosh t\le\frac{a^2z^2}2\cosh(z/2)$, using
    $\sinh v\le v\cosh v$, and
    $z^2\cosh(z/2)\le\frac{32}{\mathrm e^2}\mathrm e^{|z|}\le8\cosh z$.
\end{proof}

\begin{lemma}[Conjugation by a clipped filter]
    \label{lem:filter_conjugation}
    \lean{Entropy.sum_add_mul_norm_unitRow_le, Entropy.inner_rowWeight_inv,
        Entropy.re_sum_sub_conj_eq, Entropy.abs_re_inner_sub_inner_conj_le}
    \leanok
    \uses{lem:modular_conjugation_quadratic, lem:clipped_filter_cost}
    Let the marginal of $\Theta$ be $\operatorname{diag}(s)$ with $\sum s=1$, let
    $X=\sum_\alpha c_\alpha\otimes d_\alpha$ be Hermitian, and let $L=
        \operatorname{diag}(l)$ with $l>0$ and
    $|\log(l_j/l_k)|\le\frac a2|\log(s_j/s_k)|$ for positive $s_j,s_k$, where
    $a\le1/2$. Then
    $|\operatorname{Re}(\langle\Theta,X\Theta\rangle-\langle\Theta,(L\otimes1)
        X(L^{-1}\otimes1)\Theta\rangle)|\le4a^2\sum_\alpha\lVert c_\alpha\rVert
        \lVert d_\alpha\rVert$.
\end{lemma}
\begin{proof}\leanok
    Hermitian pairing writes the real part as
    $-\sum_{jk}(\cosh\log(l_j/l_k)-1)\operatorname{Re}T_{jk}$. For a product
    term $|T_{jk}|=\sqrt{s_js_k}\,|c_{jk}||\langle\omega_j,d\,\omega_k\rangle|$,
    Lemma~\ref{lem:clipped_filter_cost} bounds the hyperbolic coefficient
    times $\sqrt{s_js_k}$ by $2a^2(s_j+s_k)$, and the row and column estimate
    of Lemma~\ref{lem:modular_conjugation_quadratic} gives the factor
    $2\lVert c\rVert\lVert d\rVert$.
\end{proof}

\begin{lemma}[Local lifts and the trace argument]
    \label{lem:local_lift}
    \lean{Entropy.localLift, Entropy.regionState, Entropy.regionState_isHermitian,
        Entropy.localLift_mul,
        Entropy.localLift_one, Entropy.localLift_conjTranspose,
        Entropy.isSupportedOn_localLift, Entropy.IsSupportedOn.exists_localLift,
        Entropy.IsSupportedOn.mono, Entropy.IsSupportedOn.mul, Entropy.inner_localLift,
        Entropy.inner_localLift_conj_eq, Entropy.inner_conj_eq_of_isSupportedOn}
    \leanok
    \uses{lem:supported_product_decomposition}
    The lift $K\mapsto K\otimes1$ from a region $D$ is a unital multiplicative
    map compatible with adjoints, and its range is the set of operators
    supported on $D$. The regional pure-state matrix is Hermitian. Support is
    monotone in $D$ and closed under products.
    For every vector $\varphi$, $\langle\varphi,(K\otimes1)\varphi\rangle=
        \Tr\rho_{\varphi,D}K$. Consequently, if $K$ is invertible and commutes with
    $\rho_{\varphi,D}$, then
    $\langle\varphi,(K\otimes1)Y(K^{-1}\otimes1)\varphi\rangle=\langle\varphi,
        Y\varphi\rangle$ for every $Y$ supported on $D$.
\end{lemma}
\begin{proof}\leanok
    The algebraic properties hold for $K\otimes1$ and survive the relabelling.
    The regional state is the partial trace of a positive semidefinite
    matrix, and is therefore Hermitian.
    The expectation formula is the defining property of the partial trace,
    and the last identity follows from cyclicity of the trace, since
    $K^{-1}\rho_{\varphi,D}K=\rho_{\varphi,D}$.
\end{proof}

\begin{theorem}[Regional states under bijections of tensor factors]
    \label{thm:region_state_relabel}
    \lean{Entropy.siteConfigurationEquiv, Entropy.regionConfigurationEquiv,
        Entropy.complementConfigurationEquiv, Entropy.relabelVector,
        Entropy.regionState_relabel}
    \leanok
    \uses{lem:local_lift}
    Let $V,V'$ be finite sets of sites, let $e:V\to V'$ be a bijection,
    and let $n_v,n'_{v'}$ be nonnegative integer local dimensions.
    Write $E_v=\{i\in\mathbb N:i<n_v\}$ and
    $E'_{v'}=\{i\in\mathbb N:i<n'_{v'}\}$ for the local bases. Suppose that
    $f_v:E_v\to E'_{e(v)}$ is a bijection for each $v\in V$.
    These bijections induce a configuration bijection $f$ and, for
    $D\subseteq V$ and $D'=e(D)$, a regional configuration bijection $f_D$.
    Write $T\Omega(\sigma')=\Omega(f^{-1}(\sigma'))$ for the transported vector
    and $U_D$ for the permutation isometry induced by $f_D$. Then
    $\rho_{T\Omega,D'}=U_D\rho_{\Omega,D}U_D^*$.
    No normalization assumption is required.
\end{theorem}
\begin{proof}\leanok
    \uses{lem:partial_trace_petz_product_reindexing}
    Let $u',v'$ be configurations on $D'$. By
    Lemma~\ref{lem:partial_trace_petz_product_reindexing}, the bijection of
    complementary configurations changes variables in the defining partial trace:
    \begin{align}
        \rho_{T\Omega,D'}(u',v')
         & =\sum_{w'}T\Omega(u',w')\overline{T\Omega(v',w')}
        =\rho_{\Omega,D}(f_D^{-1}(u'),f_D^{-1}(v')).\notag
    \end{align}
    This is the asserted conjugation identity.
\end{proof}

\begin{theorem}[Regional entropy under bijections of tensor factors]
    \label{thm:region_entropy_relabel}
    \lean{Entropy.regionEntropy_relabel}
    \leanok
    \uses{thm:region_state_relabel, def:von_neumann_entropy}
    Under the hypotheses of Theorem~\ref{thm:region_state_relabel},
    $S_{T\Omega}(D')=S_\Omega(D)$.
\end{theorem}
\begin{proof}\leanok
    \uses{thm:region_state_relabel, thm:entropy_submatrix_equiv}
    By Theorem~\ref{thm:entropy_submatrix_equiv}, conjugation by a permutation
    isometry preserves the von Neumann entropy of the regional state.
\end{proof}


\begin{lemma}[Stationarity of one filter]
    \label{lem:filter_stationarity}
    \lean{Entropy.re_inner_commutator_eq_zero_of_isLocalMax,
        Entropy.isHermitian_of_re_trace_mul_skew_eq_zero,
        Entropy.commute_of_isHermitian_inv_mul_mul, Entropy.inner_liftProd_conj,
        Entropy.commute_regionState_of_unitary_max,
        Entropy.re_inner_commutator_eq_zero_of_isLocalExtr,
        Entropy.commute_regionState_of_unitary_extr,
        Entropy.commute_regionState_of_unitary_min}
    \leanok
    \uses{lem:local_lift}
    Let $K>0$ on a region $D$ maximize or minimize $\lVert A(K\otimes1)Cw\rVert$ over its
    unitary orbit, where $A$ is a product of invertible filters on regions
    containing $D$, each commuting with the regional state of the output
    $\psi=A(K\otimes1)Cw$ on its region. Then $K$ commutes with $\rho_{\psi,D}$.
\end{lemma}
\begin{proof}\leanok
    Differentiating along $\mathrm e^{tB}K\mathrm e^{-tB}$ with $B^*=-B$ gives
    $\operatorname{Re}\langle\psi,A(BK-KB)w'\rangle=0$. Since
    $BK-KB=(B-KBK^{-1})K$, the descending trace argument of
    Lemma~\ref{lem:local_lift} removes the outer filters, and the expectation
    of $B$ is imaginary. Hence $\operatorname{Re}\Tr(K^{-1}\rho KB)=0$ for every
    skew-Hermitian $B$, so $K^{-1}\rho K$ is Hermitian. In the eigenbasis of
    $K$ this forces $\rho K^2=K^2\rho$ and then $\rho K=K\rho$.
\end{proof}

\begin{lemma}[Simultaneous diagonalization]
    \label{lem:commuting_hermitian}
    \lean{Matrix.exists_unitary_diagonal_of_commute}
    \leanok
    Commuting Hermitian matrices are diagonal in a common orthonormal basis.
\end{lemma}
\begin{proof}\leanok
    The joint eigenspaces form an orthogonal internal direct sum; collect
    orthonormal bases of the nonzero ones.
\end{proof}

\begin{lemma}[Clipping of a stationary filter]
    \label{lem:filter_clipping}
    \lean{Entropy.deriv_nonpos_of_le_right, Entropy.eq_max_div_of_floor_optimal,
        Entropy.re_inner_deriv_nonpos_of_le_right, Entropy.hasDerivAt_conj_diagonal_rpow,
        Entropy.ratio_le_of_diagonal_max, Entropy.abs_log_rpow_sub_le_of_ratio_le}
    \leanok
    \uses{lem:filter_stationarity, lem:clipped_filter_cost}
    Let $K=U\operatorname{diag}(x^{a/2})U^*$ with $x\ge f>0$, $\sum x=1$, maximize
    the output norm among the filters $U\operatorname{diag}(y^{a/2})U^*$ with
    $y\ge f$, $\sum y=1$, and let the regional state of the output be
    $U\operatorname{diag}(q)U^*$. Then $q_k/x_k\le q_i/x_i$ whenever $x_i>f$.
    If moreover $\#\cdot f<1$, then $x_i=\max\{f,q_i/\lambda\}$ for some
    $\lambda>0$, hence $|\log(x_i/x_k)^{a/2}|\le\frac a2|\log(q_i/q_k)|$ for
    positive $q_i,q_k$.
\end{lemma}
\begin{proof}\leanok
    Moving mass from coordinate $i$ to coordinate $k$ is feasible for small
    $t>0$. Its first-order effect on the squared output norm is
    $a\operatorname{Re}\Tr\rho G$ with $G=U\operatorname{diag}(v/x)U^*$, after
    the descending trace argument, and equals $a(q_k/x_k-q_i/x_i)$. The
    floor-constrained optimality condition and the contraction of a common
    floor give the rest.
\end{proof}

\begin{theorem}[Conjugation by a filter commuting with the marginal]
    \label{thm:filter_energy_term}
    \lean{Entropy.toEuclideanLin_diagonal_kronecker_one, Entropy.inner_conj_unitary,
        Entropy.abs_re_inner_sub_conj_le_of_common_basis,
        Entropy.norm_modularExpectation_kronecker_le_of_log_clipping,
        Entropy.abs_re_inner_sub_inner_conj_le_logarithmic,
        Entropy.abs_re_inner_sub_conj_le_of_common_basis_logarithmic,
        Entropy.commute_of_isSupportedOn_disjoint}
    \leanok
    \uses{lem:filter_conjugation, lem:commuting_hermitian}
    If the marginal of a unit vector is $U\operatorname{diag}(q)U^*$ and
    $L=U\operatorname{diag}(l)U^*$ has clipped eigenvalue ratios with $a\le1/2$,
    then conjugating a Hermitian $X=\sum c_\alpha\otimes d_\alpha$ by
    $L\otimes1$ changes the real part of its expectation by at most
    $4a^2\sum\lVert c_\alpha\rVert\lVert d_\alpha\rVert$. Operators supported on
    disjoint regions commute.

    Let $\varphi$ be a unit vector whose marginal on the first factor is
    $\rho=U\operatorname{diag}(p)U^*$, with $U$ unitary, $p_i\ge0$ and
    $\sum_i p_i=1$. Let $L=U\operatorname{diag}(l)U^*$ with $l_i>0$, and
    suppose that $a\ge0$ satisfies
    \begin{align*}
        |\log l_i-\log l_k|\le\frac a2|\log p_i-\log p_k|
    \end{align*}
    whenever $p_i,p_k>0$. For Hermitian $X$ with $\|X\|\le c_0$ and a
    decomposition $X=\sum_{r=0}^{N-1} c_r\otimes d_r$ satisfying
    $\|c_r\|\|d_r\|\le c_0$, define
    \begin{align*}
        g(z)=\langle\varphi,(L^z\otimes I)X(L^{-z}\otimes I)\varphi\rangle.
    \end{align*}
    This function is entire, satisfies $|g(\mathrm iy)|\le c_0$ for real
    $y$, and obeys
    \begin{align*}
        |g(z)|\le2Nc_0
    \end{align*}
    whenever $a|\operatorname{Re}z|\le1$. If $\ell\ge1$,
    $\log(2N)\le2\ell$ and $2a\ell\le1$, then
    \begin{align*}
        \bigl|\operatorname{Re}\bigl(g(1)-g(0)\bigr)\bigr|
        \le8\mathrm e\,c_0\ell^2a^2.
    \end{align*}
    The condition $a\le1/(4\ell)$ used in
    \href{https://github.com/LionSR/TNLean/blob/67ea93e29f7c088f766b271208a2fde4cccf04a6/docs/notes/one_dimensional_note/part1_arealaw.tex#L711-L734}{Remark~6.3 of the one-dimensional note}
    is sufficient. In particular, a supported term with a decomposition of
    at most $d^2$ such products, where $d\ge1$, admits
    $\ell=\log(\mathrm e d)$.
\end{theorem}
\begin{proof}\leanok
    Transport by $U^*\otimes1$ to the Schmidt coordinates and apply
    Lemma~\ref{lem:filter_conjugation}. For disjoint supports, both operators
    are of the forms $A\otimes1$ and $1\otimes B$ across the same cut.

    \uses{lem:strip_quadratic, lem:supported_product_decomposition}
    In the common eigenbasis, write the rows of
    $(U^*\otimes I)\varphi$ as $\sqrt{p_i}\,\omega_i$. The nonzero
    $\omega_i$ are orthonormal. Clipping and
    $a|\operatorname{Re}z|\le1$ give
    \begin{align*}
        \sqrt{p_ip_k}\,
        \exp\bigl(\operatorname{Re}z(\log l_i-\log l_k)\bigr)
         & \le\sqrt{p_ip_k}\,
        \exp\bigl(\tfrac12|\log p_i-\log p_k|\bigr) \\
         & =\max\{p_i,p_k\}\le p_i+p_k
    \end{align*}
    for positive $p_i,p_k$. If either probability vanishes, its actual
    Schmidt row vanishes. The weighted row and column estimate in
    Lemma~\ref{lem:filter_conjugation} therefore bounds a product term by
    $2\|c_r\|\|d_r\|$ throughout this strip. Summing proves the asserted
    bound for $g$. On the imaginary axis $L^{\mathrm iy}$ is unitary, and
    Hermiticity gives $g(-x)=\overline{g(x)}$ for real $x$.

    If $a=0$, clipping makes the filter eigenvalues equal on the positive
    marginal support, so $g$ is constant. If $N=0$, then $X=0$. Otherwise,
    for $a>0$ put $f(w)=g(4w/a)$. On
    $|\operatorname{Re}w|\le1/4$ one has $|f(w)|\le2Nc_0$, while
    $|f(\mathrm iy)|\le c_0$. Lemma~\ref{lem:strip_quadratic}, with
    $m=2N$ and $x=a/4$, gives
    \begin{align*}
        \bigl|\operatorname{Re}\bigl(g(1)-g(0)\bigr)\bigr|
         & \le128\mathrm e\,c_0\ell^2(a/4)^2
        =8\mathrm e\,c_0\ell^2a^2.
    \end{align*}
    The smallness condition ensures $a/4\le1/(8\ell)$. Finally,
    $N\le d^2$ and $d\ge1$ imply
    $\log(2N)\le2\log(\mathrm e d)$; the supported decomposition is
    Lemma~\ref{lem:supported_product_decomposition}.
\end{proof}

\begin{proposition}[Stationary filters are clipped]
    \label{prop:stationary_filter_clipped}
    \lean{Entropy.IsFeasibleFilter, Entropy.IsFeasibleFilter.conj,
        Entropy.IsFeasibleFilter.posDef, Entropy.trace_cfc_conj_diagonal,
        Entropy.IsFeasibleFilter.of_common_basis, Entropy.trace_regionState,
        Entropy.exists_clipped_of_max_feasible}
    \leanok
    \uses{lem:filter_stationarity, lem:filter_clipping, lem:commuting_hermitian}
    A filter on $D$ is feasible at floor $f>0$ and weight $a>0$ when it is
    $U\operatorname{diag}(y^{a/2})U^*$ with $U$ unitary, $y\ge f$ and $\sum y=1$.
    Let a feasible filter $K$ maximize the output norm of
    $A(K\otimes1)Cw$ among feasible filters, where $A$ is a product of filters
    on regions containing $D$, each commuting with its regional state of the
    output $\psi\ne0$, and let $\#\cdot f<1$. Then $K$ commutes with
    $\rho_{\psi,D}$, and in a common eigenbasis $K=W\operatorname{diag}(l)W^*$,
    $\rho_{\psi,D}=W\operatorname{diag}(q)W^*$ with
    $|\log(l_i/l_k)|\le\frac a2|\log(q_i/q_k)|$ for positive $q_i,q_k$.
\end{proposition}
\begin{proof}\leanok
    Unitary conjugation preserves feasibility, so Lemma~\ref{lem:filter_stationarity}
    gives commutation. In a common eigenbasis the eigenvalues of $K$ are
    $x^{a/2}$ with $x\ge f$, $\sum x=1$, by comparing $\Tr K^{2/a}$ in the two
    bases, and diagonal variations in that basis are feasible.
    Lemma~\ref{lem:filter_clipping} applies with $\sum q=\lVert\psi\rVert^2>0$.
\end{proof}

\begin{theorem}[Positivity of feasible filters]
    \label{thm:feasible_filter_posSemidef}
    \lean{Entropy.IsFeasibleFilter.posSemidef}
    \leanok
    \uses{prop:stationary_filter_clipped}
    For real $f$ and $a$, call a matrix $K$ feasible at floor $f$ and
    weight $a$ if some unitary $U$ and weights $y_i\ge f$ with
    $\sum_i y_i=1$ satisfy
    \begin{align}
        K & =U\operatorname{diag}(y_i^{a/2})U^*.
        \label{eq:feasible_filter_coordinates}
    \end{align}
    If $f\ge0$, every such matrix is positive semidefinite, for any real
    weight $a$.
\end{theorem}
\begin{proof}
    \leanok
    Every $y_i$ is nonnegative, so every $y_i^{a/2}$ is nonnegative.
    Unitary conjugation preserves positive semidefiniteness.
\end{proof}

\begin{theorem}[Spectral decomposition and trace]
    \label{thm:hermitian_spectral_trace}
    \lean{Matrix.IsHermitian.trace_cfc_eq_sum,
        Matrix.IsHermitian.spectral_form}
    \leanok
    Let $A$ be a Hermitian matrix on a finite index set over $\mathbb R$
    or $\mathbb C$, with eigenvalues $\lambda_i$, and let
    $g:\mathbb R\to\mathbb R$ be arbitrary. Then
    $\Tr g(A)=\sum_i g(\lambda_i)$. In the complex case, if $U$ is
    its eigenvector unitary, then
    $A=U\operatorname{diag}(\lambda_i)U^*$.
\end{theorem}
\begin{proof}
    \leanok
    The spectral theorem gives the decomposition. Functional calculus
    applies $g$ to the diagonal entries; cyclicity of the trace removes
    the unitary conjugation and leaves the sum of those entries.
\end{proof}

\begin{theorem}[Trace normalization of feasible filters]
    \label{thm:feasible_filter_trace_normalization}
    \lean{Entropy.IsFeasibleFilter.trace_cfc_rpow,
        Entropy.isFeasibleFilter_zero_iff,
        Entropy.isFeasibleFilter_zero_iff_eigenvalues}
    \leanok
    \uses{prop:stationary_filter_clipped}
    Let $K$ be a complex matrix on a finite index set and let $a>0$.
    Then $K$ is feasible at floor zero and weight $a$ if and only if
    $K\ge0$ and $\Tr K^{2/a}=1$, where the power is defined by real
    continuous functional calculus. Equivalently, for $K\ge0$ with
    eigenvalues $\lambda_i$, feasibility holds if and only if
    $\sum_i\lambda_i^{2/a}=1$. The index set may be empty and $K$ may
    be singular.
\end{theorem}
\begin{proof}
    \leanok
    \uses{thm:feasible_filter_posSemidef, prop:stationary_filter_clipped,
        thm:hermitian_spectral_trace}
    For a representation in~(\ref{eq:feasible_filter_coordinates}) with
    $f=0$, positivity follows from
    Theorem~\ref{thm:feasible_filter_posSemidef}. Compute the trace in
    the basis given by $U$. Since $a>0$, the identity
    $(y_i^{a/2})^{2/a}=y_i$ holds even when $y_i=0$, so the trace is
    $\sum_i y_i=1$.

    Conversely, diagonalize $K\ge0$ and put $y_i=\lambda_i^{2/a}$.
    The trace condition gives $\sum_i y_i=1$, while $y_i\ge0$ and
    $y_i^{a/2}=\lambda_i$. This reconstructs
    (\ref{eq:feasible_filter_coordinates}) at floor zero. The same
    trace computation gives the eigenvalue formulation.
\end{proof}

The theorem identifies the matrices satisfying the trace normalization
used in the companion filter construction of
Ref.~\cite{OpenAI2026PEPSFiltersEnergy}. It does not select a unique
spectral basis. When the index set is empty, both normalized sets are
empty. The positive-weight assumption is essential: at $a=0$, the
identity matrix in dimension two is feasible, but the zeroth real power
has trace two.

\begin{theorem}[Feasibility under reindexing]
    \label{thm:feasible_filter_reindex}
    \lean{Entropy.IsFeasibleFilter.reindex,
        Entropy.isFeasibleFilter_reindex_iff}
    \leanok
    \uses{prop:stationary_filter_clipped}
    Let $e:I\to J$ be a bijection of finite index sets and let $f,a$ be
    arbitrary real numbers. Reindexing both matrix indices by $e$
    preserves feasibility at floor $f$ and weight $a$ in both directions.
\end{theorem}
\begin{proof}
    \leanok
    Transport $U$ by the same row and column bijection and replace the
    weights by $y\circ e^{-1}$. Matrix multiplication, adjoints, the
    identity, and the sum of the weights are preserved. The inverse
    bijection gives the converse.
\end{proof}

\begin{theorem}[Selected regional matrix products]
    \label{thm:selected_filter_product}
    \lean{Entropy.liftProd_chainList_eq}
    \leanok
    \uses{lem:local_lift, thm:all_filters_clipped}
    Let $(D_j)_{j\in\mathbb N}$ be regions in a finite site set, let
    $K_j$ be an arbitrary regional matrix for $j<m$, and choose a subset
    of these indices. If its elements in increasing order are
    $j_1<\cdots<j_r$, the product of the associated local lifts in the
    reversed list is $L_{j_r}\cdots L_{j_1}$. The empty list has product
    equal to the identity. No nesting or feasibility assumption is needed.
\end{theorem}
\begin{proof}
    \leanok
    Expanding the list of selected regional matrices and then taking
    their local lifts yields the same ordered list of full matrices.
\end{proof}

\begin{theorem}[Coordinate changes for matrix actions]
    \label{thm:coordinate_matrix_action}
    \lean{Matrix.reindex_mulVec, Entropy.norm_relabelVector}
    \leanok
    \uses{thm:region_state_relabel}
    Let $e:I\to J$ be a bijection of finite coordinate sets, let $R$ be
    a possibly nonunital, nonassociative semiring, let $M$ be a matrix
    over $R$ indexed by $I$, and let $v\in R^I$. Writing $M'$ for the
    matrix with both indices reindexed by $e$, one has
    $M'(v\circ e^{-1})=(Mv)\circ e^{-1}$.
    For the complex configuration vectors of
    Theorem~\ref{thm:region_state_relabel}, the corresponding map $T$
    also satisfies $\|T\Omega\|=\|\Omega\|$.
\end{theorem}
\begin{proof}
    \leanok
    Changing variables in the finite matrix-vector sum gives the first
    identity. The configuration bijection acts on Euclidean space as
    a permutation isometry, which gives the norm identity.
\end{proof}

\begin{theorem}[Successive matrix actions]
    \label{thm:ordered_matrix_action}
    \lean{Matrix.foldl_toEuclideanCLM_eq_reverse_prod,
        Entropy.foldl_localLift_eq_liftProd_chainList}
    \leanok
    \uses{lem:local_lift, thm:all_filters_clipped}
    Let $I$ be a finite set, let $\mathbb K$ be $\mathbb R$ or $\mathbb C$,
    and let $M_0,\ldots,M_{m-1}$ be matrices on $\mathbb K^I$.
    Starting at $\psi_0=\Omega$ and setting $\psi_{j+1}=M_j\psi_j$ gives
    $\psi_m=M_{m-1}\cdots M_0\Omega$. In particular, successive applications
    of the local lifts $L_j$ give their product in decreasing index order.
    No commutation or invertibility assumption is required. When $m=0$,
    the output is $\Omega$; the coordinate set $I$ may also be empty.
\end{theorem}
\begin{proof}
    \leanok
    \uses{thm:selected_filter_product}
    Induction on the list uses $\psi_{j+1}=M_j\psi_j$ and associativity of
    matrix action. Appending the first applied factor on the right of
    the remaining reversed product preserves the order of all factors.
\end{proof}

\begin{theorem}[Ordered local products under site bijections]
    \label{thm:filter_output_relabel}
    \lean{Entropy.localLift_relabel, Entropy.liftProd_chainList_relabel,
        Entropy.relabelVector_liftProd_chainList,
        Entropy.relabelVector_normalized_liftProd_chainList}
    \leanok
    \uses{lem:local_lift, thm:region_state_relabel}
    Use the site and local-basis bijections of
    Theorem~\ref{thm:region_state_relabel}. Let $(D_j)_{j\in\mathbb N}$
    and $(D'_j)_{j\in\mathbb N}$ be regions with $D'_j=e(D_j)$ for every
    $j$. For $j<m$, choose an arbitrary matrix $K_j$ on $D_j$, and let
    $K'_j=U_{D_j}K_jU_{D_j}^*$. Write $L_j$ and $L'_j$ for their local
    lifts. Then $TL_j=L'_jT$. For every selected subset of the indices
    $j<m$, let $A$ and $A'$ be the corresponding products in decreasing
    index order. For every input $\Omega$,
    \begin{align}
        TA\Omega & =A'T\Omega,\notag                   \\
        T\bigl(\|A\Omega\|^{-1}A\Omega\bigr)
                 & =\|A'T\Omega\|^{-1}A'T\Omega.\notag
    \end{align}
    Here $0^{-1}=0$. Empty regions, empty local bases, empty selected
    sets, singular matrices, and repeated regions are allowed.
\end{theorem}
\begin{proof}
    \leanok
    \uses{thm:coordinate_matrix_action, thm:selected_filter_product}
    In the region-complement coordinates, a local lift has entries
    $K_j(u,v)\delta_{w,z}$. Relabelling replaces $u,v$ by their inverse
    regional images and $w,z$ by their inverse complementary images.
    The complementary bijection preserves equality, so it preserves
    $\delta_{w,z}$. Reindexing preserves matrix multiplication and the
    identity, hence the ordered products. Changing variables in the
    matrix-vector sum proves the output identity. Since $T$ preserves
    norms and scalar multiplication, it also preserves normalization,
    even when the output vanishes.
\end{proof}

\begin{theorem}[Global comparison under site bijections]
    \label{thm:filter_max_comparison_relabel}
    \lean{Entropy.forall_norm_liftProd_chainList_le_relabel_iff}
    \leanok
    \uses{thm:filter_output_relabel, thm:feasible_filter_reindex}
    Under the hypotheses of Theorem~\ref{thm:filter_output_relabel}, fix
    arbitrary real floors $b_j$ and weights $a_j$ for $j<m$.
    Let $A(K)$ denote the product of all local lifts in decreasing index
    order. For any distinguished family $K$, the inequality
    $\|A(G)\Omega\|\le\|A(K)\Omega\|$ for every family $G$ feasible at
    $(b_j,a_j)$ holds if and only if the corresponding inequality holds
    for $K'$ and $T\Omega$ against every feasible family on the regions
    $D'_j$. The distinguished family itself need not be feasible.
\end{theorem}
\begin{proof}
    \leanok
    \uses{thm:filter_output_relabel, thm:feasible_filter_reindex}
    The regional bijections preserve feasibility and the norm of every
    ordered output. Every competitor on the target regions is the image
    of its inverse-reindexed family. Applying the inequality to that
    family proves one direction, and reindexing a source competitor
    proves the other.
\end{proof}

% These are generic matrix identities for the ascending application convention
% of OpenAI2026PEPSFiltersEnergy, Basic.lean:166--179. They do not assert equality
% with unimported source types or complete the geometric/physical comparison.

\begin{theorem}[All stationary filters are clipped]
    \label{thm:all_filters_clipped}
    \lean{Entropy.chainList, Entropy.liftProd_chainList_split,
        Entropy.chainList_update_of_ne, Entropy.isDescendingChain_map,
        Entropy.forall_clipped_of_max}
    \leanok
    \uses{prop:stationary_filter_clipped}
    Let feasible filters $K_j$ on nested regions $D_0\subseteq\dots\subseteq
        D_{m-1}$ maximize $\lVert L_{m-1}\cdots L_0\Omega\rVert$ over the product of
    the feasible sets, with nonzero output $\psi$ and $\#\cdot f_j<1$. Then for
    every $j$, $K_j$ commutes with $\rho_{\psi,D_j}$, and in a common eigenbasis
    its eigenvalue ratios are clipped.
\end{theorem}
\begin{proof}\leanok
    Descending induction on $j$: the outer filters form a descending chain by
    the induction hypothesis, and replacing $K_j$ alone keeps the product
    feasible, so Proposition~\ref{prop:stationary_filter_clipped} applies.
\end{proof}

\begin{proposition}[Existence of an optimal nested product]
    \label{prop:filter_maximizer}
    \lean{Matrix.isCompact_setOf_conjTranspose_mul_self_eq_one,
        Entropy.isCompact_floorSimplex, Entropy.isCompact_setOf_isFeasibleFilter,
        Entropy.isFeasibleFilter_uniform, Entropy.continuous_norm_liftProd_chainList,
        Entropy.exists_max_feasible_chain,
        Entropy.toEuclideanLin_liftProd_chainList_ne_zero}
    \leanok
    Suppose every region $D_j$ carries at least one configuration and
    $\#\cdot f_j\le 1$. Then the norm $\lVert L_{m-1}\cdots L_0\Omega\rVert$
    attains its maximum over families of feasible filters. If moreover all floors
    are positive and $\Omega\ne0$, every such product maps $\Omega$ to a nonzero
    vector.
\end{proposition}
\begin{proof}\leanok
    The feasible filters are the continuous image of the product of the compact
    unitary group and the compact floored simplex, and the uniform filter is
    feasible. The norm is continuous in the filters, so it attains its maximum
    on the compact product of the feasible sets. Feasible filters with positive
    floors are positive definite, hence invertible, and so is their product.
\end{proof}

\begin{lemma}[Conjugation by the splitting filter]
    \label{lem:split_filter_conjugation}
    \lean{Entropy.inner_cutOperator, Entropy.cutOperator_mul, Entropy.inv_conj_diagonal,
        Entropy.abs_re_inner_sub_conj_localLift_le}
    \leanok
    \uses{thm:filter_energy_term}
    Let $\varphi$ be a unit vector, $B$ a region, and let $L=W\operatorname{diag}(l)W^*$ and
    $\rho_{\varphi,B}=W\operatorname{diag}(q)W^*$ with $W$ unitary, $l>0$, and
    $\lvert\log l_i-\log l_k\rvert\le\tfrac a2\lvert\log q_i-\log q_k\rvert$ for positive
    $q_i,q_k$, where $a\le\tfrac12$. For a Hermitian operator $X$ supported on $S$ with
    $\lVert X\rVert\le c_0$,
    \[
        \bigl\lvert\Re\bigl(\langle\varphi,X\varphi\rangle
        -\langle\varphi,LXL^{-1}\varphi\rangle\bigr)\bigr\rvert
        \le 4a^2d_S^2c_0,\qquad d_S=\prod_{v\in S}\dim\mathcal H_v.
    \]
\end{lemma}
\begin{proof}\leanok
    In cut coordinates for $B\mid B^c$ the operator $X$ is a sum of at most $d_S^2$ products
    $c\otimes d$ with $\lVert c\rVert\lVert d\rVert\le c_0$, and
    Theorem~\ref{thm:filter_energy_term} applies.
\end{proof}

\begin{lemma}[One interaction term]
    \label{lem:chain_term_conjugation}
    \lean{Entropy.liftProdInv_append, Entropy.List.filter_not_append_filter_of_upward,
        Entropy.chainList_eq_append_of_upward, Entropy.chainList_eq_split,
        Entropy.liftProd_mul_mul_liftProdInv_of_disjoint, Entropy.splitIndices,
        Entropy.mem_splitIndices, Entropy.abs_re_inner_sub_chain_conj_le}
    \leanok
    \uses{lem:split_filter_conjugation, thm:all_filters_clipped}
    Let $K_j$ be invertible filters on nested regions $D_j$, each commuting with the regional
    state of the unit vector $\varphi$ and clipped relative to it with weight $a_j\le\tfrac12$.
    Let $X$ be Hermitian, supported on $S$, with $\lVert X\rVert\le c_0$, and suppose that at
    most one region meets $S$ without containing it. Then conjugation by the nested product
    $K$ changes $\Re\langle\varphi,X\varphi\rangle$ by at most
    $\sum_{j}4a_j^2d_S^2c_0$, the sum over the regions splitting $S$.
\end{lemma}
\begin{proof}\leanok
    The filters on regions disjoint from $S$ commute with $X$. The filters on regions
    containing $S$ drop out of the expectation by the descending trace argument. If a region
    splits $S$, the regions after it contain $S$ and the regions before it are disjoint from
    $S$, so only its conjugation remains, and Lemma~\ref{lem:split_filter_conjugation}
    applies; otherwise the expectation is unchanged.
\end{proof}

\begin{theorem}[Excitation energy of the optimal nested product]
    \label{thm:chain_energy}
    \lean{Entropy.regionState_smul, Entropy.clipped_smul, Entropy.toEuclideanLin_conj_apply_eq,
        Entropy.re_inner_sub_le_of_chain, Entropy.re_inner_sub_le_of_max}
    \leanok
    \uses{lem:chain_term_conjugation, prop:filter_maximizer}
    Let feasible filters $K_j$ on nested regions $D_j$, with floors $f_j>0$,
    $\#\cdot f_j<1$ and weights $0<a_j\le\tfrac12$, maximize
    $\lVert L_{m-1}\cdots L_0\Omega\rVert$, where $\Omega\ne0$ satisfies $H\Omega=E_0\Omega$
    for $H=\sum_ih_i$ with $h_i$ Hermitian, supported on $S_i$, $\lVert h_i\rVert\le c_0$, and
    every $S_i$ split by at most one region. Then the normalized output
    $\varphi=K\Omega/\lVert K\Omega\rVert$ satisfies
    \[
        \Re\langle\varphi,H\varphi\rangle-E_0
        \le\sum_j4a_j^2\sum_{i\text{ crossing }D_j}d_{S_i}^2c_0 .
    \]
\end{theorem}
\begin{proof}\leanok
    Since $KHK^{-1}\varphi=E_0\varphi$, the left side is the sum over $i$ of
    $\Re\bigl(\langle\varphi,h_i\varphi\rangle-\langle\varphi,Kh_iK^{-1}\varphi\rangle\bigr)$.
    By Theorem~\ref{thm:all_filters_clipped} every filter commutes with the regional state of
    $\varphi$ and is clipped relative to it; Lemma~\ref{lem:chain_term_conjugation} bounds each
    term, and exchanging the two sums gives the claim.
\end{proof}

\begin{lemma}[One step of the first norm comparison]
    \label{lem:prefix_step}
    \lean{Entropy.chainList_lt_last, Entropy.chainList_last_lt, Entropy.liftProd_chainList_succ,
        Entropy.mul_norm_le_of_overlap}
    \leanok
    Let $\psi=L_k\cdots L_0\Omega$ with $L_k$ Hermitian, suppose the normalization of $\psi$
    has overlap at least $s$ with $\Omega$, and $\lVert L_{k-1}\cdots L_0\Omega\rVert\le N$.
    Then $s\lVert\psi\rVert\le\lVert L_k\Omega\rVert N$.
\end{lemma}
\begin{proof}\leanok
    $s\lVert\psi\rVert\le\lvert\langle\Omega,\psi\rangle\rvert
        =\lvert\langle L_k\Omega,L_{k-1}\cdots L_0\Omega\rangle\rvert$, and Cauchy--Schwarz applies.
\end{proof}

\begin{lemma}[Feasible filters on a region]
    \label{lem:region_filter_norm}
    \lean{Entropy.IsFeasibleFilter.eigenvalues, Entropy.cutVector_toEuclideanLin,
        Entropy.norm_localLift_sq_le_exp}
    \leanok
    \uses{lem:filter_norm}
    Under the hypotheses of Theorem~\ref{thm:marginal_tails} for the cut $B\mid B^c$, a
    feasible filter $L$ on $B$ with positive floor and weight $0<a\le\tfrac12$, such that
    $a/(1-a)$ is admissible for the cut, satisfies
    $\lVert L\Psi\rVert^2\le\exp\bigl(-aS(\rho_{\Psi,B})+1024e\vartheta\mathcal B_Ba^2\bigr)$.
\end{lemma}
\begin{proof}\leanok
    The eigenvalues of $L$ are nonnegative with $\sum_j l_j^{2/a}=1$. In cut coordinates
    $L$ acts as $L\otimes1$, and Lemma~\ref{lem:filter_norm} applies.
\end{proof}

\begin{theorem}[The first norm comparison]
    \label{thm:first_norm_comparison}
    \lean{Entropy.chainOptimum, Entropy.chainOptimum_zero, Entropy.exists_eq_chainOptimum,
        Entropy.norm_le_chainOptimum, Entropy.pow_mul_le_prod_of_step,
        Entropy.sq_norm_inner_ge_of_gap, Entropy.sqrt_mul_chainOptimum_succ_le,
        Entropy.pow_sqrt_mul_chainOptimum_le}
    \leanok
    \uses{lem:prefix_step, thm:chain_energy, prop:filter_maximizer}
    Let $N_k$ be the largest output norm $\lVert L_{k-1}\cdots L_0\Omega\rVert$ over feasible
    filters on the first $k$ nested regions, where the floors are positive with
    $\#\cdot f_j<1$ and $0<a_j\le\tfrac12$. Let $\Omega$ be a unit vector with
    $H\Omega=E_0\Omega$ and $H-E_0\ge\Delta(1-\lvert\Omega\rangle\langle\Omega\rvert)$, where
    $\Delta>0$ and $H=\sum_ih_i$ with $h_i$ Hermitian, supported on $S_i$, $\lVert h_i\rVert\le c_0$,
    and every $S_i$ split by at most one region. If
    $\sum_{j<m}4a_j^2\sum_{i\text{ crossing }D_j}d_{S_i}^2c_0\le E$ and every feasible filter
    on $D_k$ has $\lVert L\Omega\rVert\le x_k$, then
    \[
        \sqrt{1-E/\Delta}^{\,m}N_m\le\prod_{k<m}x_k .
    \]
\end{theorem}
\begin{proof}\leanok
    The optimum is attained by Proposition~\ref{prop:filter_maximizer}. By
    Theorem~\ref{thm:chain_energy} and the gap, the normalized output of an optimal prefix of
    length $k+1$ has squared overlap at least $1-E/\Delta$ with $\Omega$, and
    Lemma~\ref{lem:prefix_step} gives $\sqrt{1-E/\Delta}\,N_{k+1}\le x_kN_k$. Telescoping from
    $N_0=1$ gives the claim.
\end{proof}

\begin{proposition}[The first norm comparison at floor zero]
    \label{prop:zero_floor_comparison}
    \lean{Entropy.isFeasibleFilter_floorRaise, Entropy.norm_le_of_forall_floor,
        Entropy.pow_sqrt_mul_norm_le_of_zero_floor}
    \leanok
    \uses{thm:first_norm_comparison}
    Under the hypotheses of Theorem~\ref{thm:first_norm_comparison} with $E<\Delta$, and bounds
    $\lVert L\Omega\rVert\le x_k$ for filters on $D_k$ feasible at any positive floor, every
    family of $m$ nonnegative filters $L_j=U_j\operatorname{diag}(y_j^{a_j/2})U_j^*$ with
    $y_j\ge0$ and $\sum y_j=1$ satisfies
    $\sqrt{1-E/\Delta}^{\,m}\lVert L_{m-1}\cdots L_0\Omega\rVert\le\prod_{k<m}x_k$.
\end{proposition}
\begin{proof}\leanok
    The filters $U_j\operatorname{diag}\bigl(((1-\delta)y_j+\delta/d_j)^{a_j/2}\bigr)U_j^*$ are
    feasible at the floor $\delta/d_j$ and converge to $L_j$ as $\delta\downarrow0$. The bound
    of Theorem~\ref{thm:first_norm_comparison} does not depend on the floor, and the output
    norm is continuous in the filters.
\end{proof}


\begin{theorem}[One logarithmic prefix comparison]
    \label{thm:logarithmic_prefix_step}
    \lean{Entropy.sqrt_mul_chainOptimum_succ_le_logarithmic}
    \leanok
    Let $\mathcal H=\bigotimes_{v\in V}\mathbb C^{n_v}$ be a finite
    tensor product, let $D_0\subseteq D_1\subseteq\cdots$, and let
    $a_j>0$ and $f_j>0$ satisfy $f_j\dim\mathcal H_{D_j}<1$.
    A feasible filter on $D_j$ has the form
    $U\operatorname{diag}(y_r^{a_j/2})U^*$, where $U$ is unitary,
    $y_r\ge f_j$ and $\sum_r y_r=1$. Let $N_r$ be the maximum of
    $\lVert\widehat L_{r-1}\cdots\widehat L_0\Omega\rVert$ over
    these feasible families, where $\widehat L_j=L_j\otimes I_{D_j^c}$.
    Suppose that $\Omega$ is a unit vector, $H=\sum_{i\in I}h_i$,
    and each $h_i$ is Hermitian, supported on $S_i$, with
    $\lVert h_i\rVert\le c_0$. Assume
    \begin{align*}
        H\Omega & =E_0\Omega,                                          \\
        H-E_0I  & \ge\Delta(I-\lvert\Omega\rangle\langle\Omega\rvert),
        \qquad \Delta>0.
    \end{align*}
    Each $S_i$ is split by at most one of the regions. Put
    \begin{align*}
        \ell_i       & =\log\Bigl(\mathrm e\prod_{v\in S_i}n_v\Bigr), \\
        \mathcal C_j & =\{i\in I:S_i\cap D_j\ne\varnothing,
        \ S_i\setminus D_j\ne\varnothing\},                           \\
        B_r          & =8\mathrm e\sum_{j=0}^{r-1}a_j^2
        \sum_{i\in\mathcal C_j}c_0\ell_i^2.
    \end{align*}
    Fix $k\ge0$ and $E,x\in\mathbb R$. Assume $B_{k+1}\le E$,
    and assume $2a_j\ell_i\le1$ for every $j<k+1$ and
    $i\in\mathcal C_j$. If every feasible filter $L$ on $D_k$
    satisfies $\lVert\widehat L\Omega\rVert\le x$, then
    \begin{align*}
        \sqrt{1-E/\Delta}\,N_{k+1}\le xN_k.
    \end{align*}
    No bound on $E$ relative to $\Delta$ is required. The square
    root is the nonnegative real square root, equal to zero for a
    nonpositive argument. This is the prefix comparison in
    \cite[Lemma~3.2]{OpenAI2026AreaLaw}, with the local energy
    coefficient from
    \href{https://github.com/LionSR/TNLean/blob/67ea93e29f7c088f766b271208a2fde4cccf04a6/docs/notes/one_dimensional_note/part1_arealaw.tex#L711-L734}{Remark~6.3 of the one-dimensional note}.
\end{theorem}
\begin{proof}
    \leanok
    \uses{thm:filter_max_logarithmic_energy, lem:prefix_step,
        prop:filter_maximizer}
    Choose an actual maximizing family of length $k+1$. Positive
    floors make its output $\psi$ nonzero. The logarithmic energy
    estimate applies to the normalization of this same $\psi$,
    and the gap gives squared overlap at least $1-E/\Delta$.
    Monotonicity of the real square root gives overlap at least
    $\sqrt{1-E/\Delta}$. Splitting off the last filter and using
    Lemma~\ref{lem:prefix_step} bounds the preceding output by
    $N_k$ and the last single-filter norm by $x$. No independent
    normalized output or clipping data are assumed.
\end{proof}

\begin{theorem}[Iterated logarithmic prefix comparison]
    \label{thm:logarithmic_prefix_product}
    \lean{Entropy.pow_sqrt_mul_chainOptimum_le_logarithmic}
    \leanok
    Keep the system, positive floors, positive weights and
    Hamiltonian hypotheses of
    Theorem~\ref{thm:logarithmic_prefix_step}. Fix $m\ge0$ and
    $E\in\mathbb R$. Assume $B_m\le E$ and $2a_j\ell_i\le1$
    for every $j<m$ and $i\in\mathcal C_j$. Let $x_j\ge0$,
    and suppose that for $j<m$ every feasible filter on $D_j$
    satisfies $\lVert\widehat L\Omega\rVert\le x_j$. Then
    \begin{align*}
        \bigl(\sqrt{1-E/\Delta}\bigr)^mN_m
        \le\prod_{j=0}^{m-1}x_j.
    \end{align*}
    No additional sign or upper bound on $E$ is assumed.
\end{theorem}
\begin{proof}
    \leanok
    \uses{thm:logarithmic_prefix_step, thm:first_norm_comparison}
    Every summand of $B_m$ is nonnegative: for each actual
    interaction, $\lVert h_i\rVert\le c_0$ implies $c_0\ge0$.
    Thus $B_{k+1}\le B_m$ for $k<m$. The same $E$ and the
    restricted crossing-pair smallness conditions apply to every
    prefix. Iterate Theorem~\ref{thm:logarithmic_prefix_step},
    using $x_j\ge0$ and $N_0=\lVert\Omega\rVert=1$.
    Empty interaction families and empty products require no
    separate nonemptiness assumption.
\end{proof}

\begin{theorem}[Logarithmic norm comparison at floor zero]
    \label{thm:logarithmic_zero_floor_comparison}
    \lean{Entropy.pow_sqrt_mul_norm_le_of_zero_floor_logarithmic}
    \leanok
    Keep the regions, positive weights and Hamiltonian hypotheses
    of Theorem~\ref{thm:logarithmic_prefix_step}, without fixing
    positive floors. Fix $m\ge0$ and $E<\Delta$, and assume
    $B_m\le E$ and $2a_j\ell_i\le1$ for every $j<m$ and
    $i\in\mathcal C_j$. Let $x_j\ge0$. For $j<m$, suppose that
    $\lVert\widehat L\Omega\rVert\le x_j$ for every positive
    floor $f$ and every filter feasible at that floor and weight
    $a_j$. Every family feasible at floor zero, namely
    \begin{align*}
        K_j     & =U_j\operatorname{diag}(y_{j,r}^{a_j/2})U_j^*, \\
        y_{j,r} & \ge0,\qquad \sum_r y_{j,r}=1,
    \end{align*}
    then satisfies
    \begin{align*}
        \bigl(\sqrt{1-E/\Delta}\bigr)^m
        \lVert\widehat K_{m-1}\cdots\widehat K_0\Omega\rVert
        \le\prod_{j=0}^{m-1}x_j.
    \end{align*}
    Both $E$ and the bounds $x_j$ are uniform in the positive
    floors. The floor-zero filters may be singular; no clipping
    condition is imposed on them.
\end{theorem}
\begin{proof}
    \leanok
    \uses{thm:logarithmic_prefix_product, prop:zero_floor_comparison}
    Write $d_j=\dim\mathcal H_{D_j}$, which is positive since
    $\Omega\ne0$. For $0<\delta<1$, replace the spectral parameters by
    $(1-\delta)y_{j,r}+\delta/d_j$. The resulting filters are
    feasible at floors $\delta/d_j$ and converge to $K_j$.
    The energy budget and smallness conditions contain no floor,
    and the single-filter bounds are uniform in that floor.
    Apply Theorem~\ref{thm:logarithmic_prefix_product} at each
    raised floor. Since $E<\Delta$, the square-root factor is
    positive; the uniform bound passes to the floor-zero output
    by continuity, as in Proposition~\ref{prop:zero_floor_comparison}.
\end{proof}

\begin{theorem}[Closed conditions for maximizing filters]
    \label{thm:filter_zero_floor_closed_max}
    \lean{Entropy.exists_max_feasible_chain_mem_of_pos_floor}
    \leanok
    \uses{prop:stationary_filter_clipped}
    Let $D_0,\ldots,D_{m-1}$ be regions of a finite system, with
    nonempty regional configuration spaces and weights $a_j\ge0$.
    Fix an input vector $\Omega$. Suppose a closed set of filter families
    contains a norm-maximizing family at every floor $\delta/d_j$,
    where $0<\delta<1$ and $d_j=\dim\mathcal H_{D_j}$.
    Then it contains a maximizing family at floor zero.
\end{theorem}
\begin{proof}\leanok
    \uses{prop:filter_maximizer,prop:zero_floor_comparison,thm:buffer_discount}
    Positive-floor feasible families belong to the compact set of
    zero-floor feasible families. Choose $\delta_k=1/(2(k+1))$ and a
    convergent subsequence of the corresponding maximizers. Closedness
    preserves the stated condition. Every family at a fixed positive
    floor is admissible along a sufficiently late part of the subsequence,
    so continuity bounds its output norm by the limiting norm.
    Approximation by positive-floor families gives the same comparison
    at floor zero.
\end{proof}

\begin{lemma}[Nonzero output from norm dominance]
    \label{lem:filter_zero_floor_nonzero}
    \lean{Entropy.norm_liftProd_pos_of_zero_floor_max}
    \leanok
    \uses{prop:stationary_filter_clipped}
    Let $D_0,\ldots,D_{m-1}$ be regions of a finite system with nonempty
    regional configuration spaces, and let $a_j$ be arbitrary real weights.
    If $\Omega\ne0$ and a filter family has output norm at least that of
    every floor-zero feasible family, then its output is nonzero.
    The family being compared need not itself be feasible.
\end{lemma}
\begin{proof}\leanok
    \uses{prop:filter_maximizer,thm:buffer_discount}
    The uniform family is feasible at floors $1/(2d_j)$, where
    $d_j=\dim\mathcal H_{D_j}$, and its product is invertible.
    Its output on $\Omega$ has positive norm, which is a lower bound
    for the output norm of the given family.
\end{proof}

\begin{theorem}[Excitation energy without a positive floor]
    \label{thm:filter_zero_floor_energy}
    \lean{Entropy.exists_max_feasible_chain_zero_floor_energy,
        Entropy.exists_max_feasible_chain_zero_floor_normalized_energy}
    \leanok
    \uses{lem:local_lift,prop:stationary_filter_clipped}
    Let $D_0\subseteq\cdots\subseteq D_{m-1}$ be regions of a finite system,
    and let $0<a_j\le1/2$. Suppose $\Omega\ne0$ and $H\Omega=E_0\Omega$,
    where $H=\sum_i h_i$, each $h_i$ is Hermitian, supported on $S_i$,
    and $\lVert h_i\rVert\le c_0$. Assume each $S_i$ is split by at most
    one $D_j$. Put
    \begin{align}
        B & =\sum_j4a_j^2\sum_{i:\,D_j\text{ splits }S_i}
        (\dim\mathcal H_{S_i})^2c_0.
    \end{align}
    There is a floor-zero feasible family that maximizes the output norm
    and has nonzero output $\psi=K_{m-1}\cdots K_0\Omega$ satisfying
    \begin{align}
        \operatorname{Re}\langle\psi,H\psi\rangle-E_0\lVert\psi\rVert^2
         & \le B\lVert\psi\rVert^2.
        \label{eq:filter_zero_floor_energy}
    \end{align}
    Thus $\phi=\psi/\lVert\psi\rVert$ has unit norm and
    $\operatorname{Re}\langle\phi,H\phi\rangle-E_0\le B$.
    The maximizing filters may be singular.
\end{theorem}
\begin{proof}\leanok
    \uses{thm:chain_energy,thm:filter_zero_floor_closed_max,
        lem:filter_zero_floor_nonzero,prop:filter_maximizer}
    At every positive floor, the maximizing-family energy estimate
    implies (\ref{eq:filter_zero_floor_energy}). Its two sides are
    continuous functions of the filter matrices, so it defines a closed
    set even at zero output. Theorem~\ref{thm:filter_zero_floor_closed_max}
    supplies a maximizing family in that set.
    Lemma~\ref{lem:filter_zero_floor_nonzero} makes its output nonzero.
    Divide by its squared norm.
\end{proof}

\begin{lemma}[Trial filters on a disjoint union]
    \label{lem:trial_filters}
    \lean{Entropy.localLift_apply, Entropy.regionUnionEquiv, Entropy.localLift_union_kronecker,
        Entropy.isFeasibleFilter_union_kronecker, Entropy.pinProj, Entropy.pinProj_conjTranspose,
        Entropy.pinProj_mul_self, Entropy.sum_pinProj, Entropy.eq_of_forall_trace_mul_eq,
        Entropy.norm_toEuclideanLin_sq, Entropy.norm_sq_localLift_pinProj,
        Entropy.inner_localLift_eq_sum_pinProj, Entropy.regionState_eq_sum_pinProj,
        Entropy.mul_prod_map_mul_of_idempotent, Entropy.mul_prod_eq_prod_mul,
        Entropy.norm_toEuclideanLin_le_of_proj, Entropy.norm_liftProd_trial_ge}
    \leanok
    Let $\rho_{\Omega,X}=U\operatorname{diag}(p)U^*$ and let $P_i$ be the projection onto the
    $i$-th eigenvector $u_i$. Then $\lVert P_i\Omega\rVert^2=p_i$, and on every region $T$
    disjoint from $X$ the regional state of $\Omega$ is the sum of those of the vectors
    $P_i\Omega$. For regions $T_j$ disjoint from $X$ and filters $M_j$ on $T_j$, the filters
    $P_i\otimes M_j$ on $X\cup T_j$ are nonnegative with $\operatorname{Tr}(P_i\otimes
        M_j)^{2/a}=\operatorname{Tr}M_j^{2/a}$, and
    $\lVert M_{m-1}\cdots M_0P_i\Omega\rVert\le\lVert(P_i\otimes M_{m-1})\cdots
        (P_i\otimes M_0)\Omega\rVert$.
\end{lemma}
\begin{proof}\leanok
    On a disjoint union the lift of $A\otimes B$ is the product of the lifts of $A$ and of $B$.
    The projection $P_i$ commutes with the filters on the $T_j$, is idempotent, and does not
    increase norms.
\end{proof}

\begin{theorem}[Complex powers and the three-lines bound]
    \label{thm:filter_three_lines}
    \lean{Entropy.specPow, Entropy.specPow_zero, Entropy.specPow_one, Entropy.specPow_add,
        Entropy.specPow_conjTranspose, Entropy.specPow_conjTranspose_mul_self,
        Entropy.hasDerivAt_specPow, Entropy.differentiable_specPow, Entropy.norm_specPow_le,
        Entropy.localLift_injective, Entropy.localLiftCLM, Entropy.chainPow,
        Entropy.chainPow_zero_length, Entropy.chainPow_succ, Entropy.chainPow_zero,
        Entropy.chainPow_conjTranspose_mul_self, Entropy.chainPow_mul_conjTranspose_self,
        Entropy.differentiable_chainPow, Entropy.norm_chainPow_le,
        Entropy.hasDerivAt_chainPow_zero, Entropy.chainPow_eq_liftProd, Entropy.chainPow_one_add,
        Entropy.applyCLM, Entropy.norm_toEuclideanLin_chainPow_le_rpow,
        Entropy.sum_re_inner_log_le}
    \leanok
    Let $M_j=W_j\operatorname{diag}(l_j)W_j^*$ with $l_j>0$ be filters on nested regions $T_j$,
    let $\xi$ be a unit vector, and suppose that every nested product of unitary conjugates
    $V_j^*M_jV_j$, with $V_j$ unitary on $T_j$, has output norm at most $B>0$ on $\xi$. Then
    $\lVert M_{k-1}^x\cdots M_0^x\xi\rVert\le B^x$ for $0\le x\le1$, and
    $\sum_{j<k}\Re\langle\xi,\log M_j\,\xi\rangle\le\log B$.
\end{theorem}
\begin{proof}\leanok
    The function $F(z)=M_{k-1}^z\cdots M_0^z\xi$ is entire and bounded on the strip
    $0\le\Re z\le1$. On $\Re z=0$ every factor is unitary. On $\Re z=1$ the identity
    $M_{k-1}^{1+iy}\cdots M_0^{1+iy}=(U_{k-1}\cdots U_0)(V_{k-2}^*M_{k-1}V_{k-2})\cdots M_0$,
    with $U_j=M_j^{iy}$ and $V_j=U_j\cdots U_0$, writes $F(1+iy)$ as a unitary applied to a
    nested product of unitary conjugates. The three-lines theorem gives $\lVert F(x)\rVert\le
        B^x$. Since $F(0)=\xi$, the one-sided derivative of $\lVert F(x)\rVert^2-B^{2x}$ at $x=0$
    is nonpositive, which is the second claim.
\end{proof}

\begin{theorem}[Strong subadditivity for regions]
    \label{thm:region_ssa}
    \lean{Entropy.regionEntropy, Entropy.trace_regionState_submatrix_mul,
        Entropy.disjoint_union_of_disjoint, Entropy.tripleEquiv, Entropy.tripleEquiv_apply,
        Entropy.localLift_tripleEquiv_kronecker_one, Entropy.regionState_posSemidef,
        Entropy.tripleState, Entropy.partialTraceLeft_tripleState,
        Entropy.partialTraceRight_regionState_union,
        Entropy.partialTraceRight_reassoc_tripleState, Entropy.regionEntropy_ssa,
        Entropy.regionEntropy_cond_anti}
    \leanok
    For pairwise disjoint regions $A,B,C$,
    $S(A\cup B\cup C)+S(B)\le S(A\cup B)+S(B\cup C)$. In particular
    $S(X\mid T)\le S(X\mid T')$ for $T'\subseteq T$ disjoint from $X$.
\end{theorem}
\begin{proof}\leanok
    The regional states of $A\cup B$, $B\cup C$ and $B$ are the partial traces of the regional
    state of $A\cup B\cup C$, viewed as a tripartite matrix, so strong subadditivity applies.
\end{proof}

\begin{theorem}[The entropy discount from a buffer]
    \label{thm:buffer_discount}
    \lean{Entropy.IsFeasibleFilter.mono_floor, Entropy.exists_trial_filter,
        Entropy.neg_sum_le_log_of_trial, Entropy.sum_norm_sq_mul_regionEntropy_le,
        Entropy.neg_two_mul_log_le_entropy, Entropy.sum_mul_condEntropy_le,
        Entropy.two_mul_condEntropy_le}
    \leanok
    \uses{prop:zero_floor_comparison, lem:region_filter_norm, lem:trial_filters,
        thm:filter_three_lines, thm:region_ssa}
    Let $X$ be a region, $T_0\subseteq T_1\subseteq\cdots$ regions disjoint from $X$, and
    $D_j=X\cup T_j$, with weights $0<a_j\le\tfrac12$ of total mass
    $\sum_{j<m}a_j=2$. Let $\Omega$, $H$ and $\Delta$ be as in
    Theorem~\ref{thm:first_norm_comparison} for the regions $D_j$, with energy budget at most
    $E<\Delta$, and suppose every $a_j/(1-a_j)$ is admissible for the cut $D_j$. If every $T_j$
    lies in a region $T$ disjoint from $X$, then
    \[
        2S(X\mid T)\le S(X)+\sum_{j<m}1024e\vartheta\mathcal B_ja_j^2
        -m\log(1-E/\Delta).
    \]
\end{theorem}
\begin{proof}\leanok
    By Proposition~\ref{prop:zero_floor_comparison} and Lemma~\ref{lem:region_filter_norm}, every
    nested product of nonnegative filters on the $D_j$ has output norm at most
    $B=\prod_k x_k/\sqrt{1-E/\Delta}^{\,m}$ with $x_k^2=\exp(-a_kS(D_k)+1024e\vartheta\mathcal
        B_ka_k^2)$. For a Schmidt index $i$ with $p_i>0$, Lemma~\ref{lem:trial_filters} bounds the
    nested products of filters on the $T_j$ applied to $\xi_i=P_i\Omega/\sqrt{p_i}$ by
    $B/\sqrt{p_i}$. Theorem~\ref{thm:filter_three_lines} with the filters
    $((1-\delta)\rho_{\xi_i,T_j}+\delta/d)^{a_j/2}$ gives, as $\delta\downarrow0$,
    $-\sum_j\tfrac{a_j}2S_{\xi_i}(T_j)\le\log B-\tfrac12\log p_i$. Averaging with weights $p_i$
    and using concavity of the entropy gives $-2\log B\le S(X)+\sum_ja_jS(T_j)$. Hence
    $\sum_ja_jS(X\mid T_j)\le S(X)+\text{error}$, and Theorem~\ref{thm:region_ssa} bounds the
    left side below by $2S(X\mid T)$.
\end{proof}

\begin{theorem}[Weighted conditional entropy with logarithmic energy cost]
    \label{thm:buffer_weighted_logarithmic}
    \lean{Entropy.sum_mul_condEntropy_le_logarithmic}
    \leanok
    Let $\mathcal H=\bigotimes_{v\in V}\mathbb C^{n_v}$ be a finite
    tensor product. Let $X$ be a region and let
    $T_0\subseteq T_1\subseteq\cdots$ be regions disjoint from $X$.
    Put $D_j=X\cup T_j$, and let $0<a_j\le\tfrac12$.
    Suppose that $\Omega$ is a unit vector and $H=\sum_{i\in I}h_i$,
    where each $h_i$ is Hermitian, supported on $S_i$, and
    $\lVert h_i\rVert\le c_0$, with $c_0\ge0$. Assume
    \begin{align*}
        H\Omega & =E_0\Omega,                                          \\
        H-E_0I  & \ge\Delta(I-\lvert\Omega\rangle\langle\Omega\rvert),
        \qquad \Delta>0.
    \end{align*}
    Each $S_i$ is split by at most one $D_j$. Define
    \begin{align*}
        d_i          & =\prod_{v\in S_i}n_v,\qquad \ell_i=\log(\mathrm e d_i), \\
        \mathcal C_j & =\{i\in I:S_i\cap D_j\ne\varnothing,
        \ S_i\setminus D_j\ne\varnothing\},                                    \\
        \mathcal B_j & =1+\sum_{i\in\mathcal C_j}\ell_i^2,
        \qquad \vartheta=c_0/\Delta,                                           \\
        \mathcal E_m & =8\mathrm e\sum_{j=0}^{m-1}a_j^2
        \sum_{i\in\mathcal C_j}c_0\ell_i^2.
    \end{align*}
    Fix $m\ge0$ and $E<\Delta$. Suppose that $\mathcal E_m\le E$
    and, for every $j<m$, that
    \begin{align*}
        \frac{a_j}{1-a_j}
        \le\frac{1}{32\sqrt{(1+\vartheta)\mathcal B_j}}.
    \end{align*}
    Writing $S(X\mid T)=S(X\cup T)-S(T)$, one has
    \begin{align*}
        \sum_{j=0}^{m-1}a_jS(X\mid T_j)
        \le S(X)+1024\mathrm e\vartheta
        \sum_{j=0}^{m-1}\mathcal B_j a_j^2
        -m\log(1-E/\Delta).
    \end{align*}
    This is the weighted norm-comparison argument in
    \cite[Lemma~3.2]{OpenAI2026AreaLaw}, with the logarithmic energy
    coefficient of
    \href{https://github.com/LionSR/TNLean/blob/67ea93e29f7c088f766b271208a2fde4cccf04a6/docs/notes/one_dimensional_note/part1_arealaw.tex#L711-L734}{Remark~6.3 of the one-dimensional note}.
\end{theorem}
\begin{proof}
    \leanok
    \uses{thm:logarithmic_zero_floor_comparison, lem:region_filter_norm,
        lem:trial_filters, thm:filter_three_lines, thm:buffer_discount}
    The admissibility assumption and $\ell_i^2\le\mathcal B_j$
    for $i\in\mathcal C_j$ give
    $32a_j\ell_i\le1-a_j$, hence $2a_j\ell_i\le1$.
    Thus no further smallness assumption is needed to apply
    Theorem~\ref{thm:logarithmic_zero_floor_comparison}.
    Lemma~\ref{lem:region_filter_norm} bounds every single-filter
    norm by $x_j$, where
    \begin{align*}
        x_j^2=\exp\bigl(-a_jS(D_j)
        +1024\mathrm e\vartheta\mathcal B_j a_j^2\bigr).
    \end{align*}
    Consequently every feasible product at floor zero has norm at most
    $B=\prod_{j<m}x_j/(\sqrt{1-E/\Delta})^m$.
    Schmidt pinning, the trial filters and the three-lines argument
    in Theorem~\ref{thm:buffer_discount} give
    \begin{align*}
        -2\log B & \le S(X)+\sum_{j=0}^{m-1}a_jS(T_j), \\
        2\log B  & =-\sum_{j=0}^{m-1}a_jS(D_j)
        +1024\mathrm e\vartheta\sum_{j=0}^{m-1}\mathcal B_j a_j^2
        -m\log(1-E/\Delta).
    \end{align*}
    Combining these inequalities proves the assertion. The same
    $E$ and $x_j$ apply at every positive floor before passage to
    floor zero. No total-mass condition on the weights is used.
\end{proof}

\begin{theorem}[The entropy discount from a buffer with logarithmic energy cost]
    \label{thm:buffer_discount_logarithmic}
    \lean{Entropy.two_mul_condEntropy_le_logarithmic}
    \leanok
    Under the hypotheses of
    Theorem~\ref{thm:buffer_weighted_logarithmic}, assume in addition
    that $\sum_{j<m}a_j=2$. Let $T'$ be disjoint from $X$ and
    contain every $T_j$ with $j<m$. Then
    \begin{align*}
        2S(X\mid T')\le S(X)
        +1024\mathrm e\vartheta\sum_{j=0}^{m-1}\mathcal B_j a_j^2
        -m\log(1-E/\Delta).
    \end{align*}
    This is the conditional-entropy conclusion in
    \cite[Lemma~3.2]{OpenAI2026AreaLaw}, with the same logarithmic
    energy budget $\mathcal E_m$.
\end{theorem}
\begin{proof}
    \leanok
    \uses{thm:buffer_weighted_logarithmic, thm:region_ssa}
    Strong subadditivity gives $S(X\mid T')\le S(X\mid T_j)$
    for every $j<m$. Multiply by $a_j>0$, sum, and use the
    total mass two. The weighted estimate of
    Theorem~\ref{thm:buffer_weighted_logarithmic} completes the proof.
\end{proof}

\subsection{Sharp one-crossing estimate for polynomial patches}

The polynomial-PEPS patch argument uses the explicit constant in
\cite{OpenAI2026PolynomialPEPS}, Section~4,
\texttt{03-patches.tex}, lines 248--319. Its sharper estimate follows
from the same Hermitian pairing and weighted row and column bounds.

\begin{theorem}[Quadratic hyperbolic scaling]
    \label{lem:patch_cosh_scaling}
    \lean{Entropy.cosh_mul_sub_one_le_sq_mul,
        Entropy.sqrt_mul_cosh_sub_one_le_sharp}
    \leanok
    For $0\le a\le1$ and $y\in\R$,
    \begin{align*}
        \cosh(ay)-1\le a^2(\cosh y-1).
    \end{align*}
    If $p,q\ge0$, $a\le1$, and
    $|t|\le\frac a2|\log p-\log q|$ whenever $p,q>0$, then
    \begin{align*}
        \sqrt{pq}(\cosh t-1)\le\frac{a^2}{2}(p+q).
    \end{align*}
\end{theorem}
\begin{proof}\leanok
    In the power series after its constant term, every exponent is at
    least two, so $a^{2k}\le a^2$ for $k\ge1$. For positive $p,q$, set
    $z=(\log p-\log q)/2$. The clipping condition and monotonicity of
    $\cosh$ in the absolute value of its argument give the bound, using
    $\sqrt{pq}\cosh z=(p+q)/2$. If $a<0$, the clipping condition forces
    $t=0$. If either probability vanishes, the left side vanishes.
\end{proof}

\begin{theorem}[Inside-support matrix-unit count]
    \label{lem:patch_inside_decomposition}
    \lean{Entropy.card_innerConfig, Entropy.IsSupportedOn.hasProductDecomposition_inner}
    \leanok
    Suppose the physical configuration space is nonempty. An operator
    $X$ supported on $S$ has a product decomposition across $B\sqcup B^c$
    with at most $d_{S\cap B}^2$ terms, each of product norm at most
    $\lVert X\rVert$, where $d_{S\cap B}=\prod_{v\in S\cap B}n_v$.
\end{theorem}
\begin{proof}\leanok
    \uses{lem:supported_product_decomposition}
    Keep the actual number of matrix units on $S\cap B$ in the
    construction of Lemma~\ref{lem:supported_product_decomposition}.
    No enlargement to the dimension of all of $S$ is needed.
\end{proof}

\begin{theorem}[Sharp clipped conjugation]
    \label{thm:patch_sharp_conjugation}
    \lean{Entropy.abs_re_inner_sub_inner_conj_le_sharp,
        Entropy.abs_re_inner_sub_conj_le_of_common_basis_sharp,
        Entropy.abs_re_inner_sub_conj_localLift_le_sharp}
    \leanok
    Let $\phi$ be a unit vector. Suppose the positive invertible filter
    $L$ and its marginal have a common orthonormal eigenbasis, with
    eigenvalues $l_i>0$ and $p_i\ge0$. Assume $a\le1$ and
    $|\log l_i-\log l_k|\le\frac a2|\log p_i-\log p_k|$ whenever
    $p_i,p_k>0$. For a Hermitian product sum $X=\sum_\alpha c_\alpha\otimes d_\alpha$,
    \begin{align*}
        |\operatorname{Re}\langle\phi,X\phi\rangle
        -\operatorname{Re}\langle\phi,(L\otimes1)X(L^{-1}\otimes1)\phi\rangle|
        \le a^2\sum_\alpha\lVert c_\alpha\rVert\lVert d_\alpha\rVert.
    \end{align*}
    If $X$ is supported on $S$ and $\lVert X\rVert\le J$, the expectation
    difference is at most $a^2d_{S\cap B}^2J$, by choosing the inside-support
    matrix-unit decomposition.
\end{theorem}
\begin{proof}\leanok
    \uses{lem:patch_cosh_scaling, lem:filter_conjugation,
        lem:patch_inside_decomposition}
    Use the coefficient $a^2(p_i+p_k)/2$ in the Hermitian pairing.
    The weighted row and column estimate contributes the factor two.
    Unitary transport gives the common-basis statement. For supported
    operators use the matrix-unit decomposition just proved.
\end{proof}

\begin{theorem}[A single crossing in a nested chain]
    \label{thm:patch_sharp_chain}
    \lean{Entropy.abs_re_inner_sub_chain_conj_le_sharp}
    \leanok
    Let $K_j$ be invertible filters on nested regions $D_j$. Each filter
    commutes with the regional marginal of the same unit vector $\phi$
    and satisfies the common-basis and clipping hypotheses above, with
    weight $a_j\le1$. Suppose the support $S$ of a Hermitian operator $X$
    is split by at most one region, and $\lVert X\rVert\le J$. Writing
    $K=K_{m-1}\cdots K_0$ for the lifted product,
    \begin{align*}
        |\operatorname{Re}\langle\phi,X\phi\rangle
        -\operatorname{Re}\langle\phi,KXK^{-1}\phi\rangle|
        \le\sum_{j\text{ splits }S}a_j^2d_{S\cap D_j}^2J.
    \end{align*}
    In particular, an interaction with one inside endpoint of dimension
    $q$ has one-term cost at most $q^2Ja_j^2$.
\end{theorem}
\begin{proof}\leanok
    \uses{thm:patch_sharp_conjugation}
    Filters outside the crossing region drop out by the descending
    trace identity; filters inside it have disjoint support and commute
    with $X$. Apply the sharp supported bound to the remaining filter.
    If there is no crossing, the expectation difference is zero.
\end{proof}

\begin{theorem}[Logarithmic cost of one splitting filter]
    \label{thm:split_filter_conjugation_logarithmic}
    \lean{Entropy.abs_re_inner_sub_conj_localLift_le_logarithmic}
    \leanok
    Let $\varphi$ be a unit vector in a finite tensor product
    $\mathcal H=\bigotimes_{v\in V}\mathbb C^{n_v}$, let $B,S\subseteq V$,
    and put
    \begin{align*}
        d_S=\prod_{v\in S}n_v,\qquad \ell_S=\log(\mathrm e d_S).
    \end{align*}
    Suppose that the positive invertible filter $L$ on $\mathcal H_B$
    and the actual marginal $\rho_{\varphi,B}$ have a common
    orthonormal eigenbasis, with eigenvalues $l_i>0$ and $p_i\ge0$.
    Let $a\ge0$. Assume that for every pair $i,k$ with
    $p_i,p_k>0$,
    \begin{align*}
        |\log l_i-\log l_k|
        \le\frac a2|\log p_i-\log p_k|.
    \end{align*}
    Assume also that
    \begin{align*}
        2a\ell_S\le1.
    \end{align*}
    For every Hermitian operator $X$ supported on $S$ with
    $\lVert X\rVert\le c_0$, writing
    $\widehat L=L\otimes I_{B^c}$, one has
    \begin{align*}
        \left|\operatorname{Re}\left(
        \langle\varphi,X\varphi\rangle
        -\langle\varphi,\widehat L X\widehat L^{-1}\varphi\rangle
        \right)\right|
        \le8\mathrm e\,c_0\ell_S^2a^2.
    \end{align*}
    The smallness condition in
    \href{https://github.com/LionSR/TNLean/blob/67ea93e29f7c088f766b271208a2fde4cccf04a6/docs/notes/one_dimensional_note/part1_arealaw.tex#L711-L734}{Remark~6.3 of the one-dimensional note}
    is sufficient for the displayed hypothesis.
\end{theorem}
\begin{proof}\leanok
    \uses{thm:filter_energy_term, lem:patch_inside_decomposition}
    The existence of the unit vector implies $n_v\ge1$ at every site,
    so $d_S\ge1$ and $\ell_S\ge1$. By
    Theorem~\ref{lem:patch_inside_decomposition}, the operator in cut
    coordinates has a product decomposition with
    $N\le d_{S\cap B}^2\le d_S^2$ terms and product norms at most
    $c_0$. Hence
    \begin{align*}
        \log(2N)\le2\log(\mathrm e d_S)=2\ell_S;
    \end{align*}
    if $N=0$, the operator vanishes and the conclusion is immediate.
    Relabelling tensor coordinates preserves the operator norm and
    the expectation. The logarithmic estimate of
    Theorem~\ref{thm:filter_energy_term}, applied to the common
    eigenbasis of the filter and the actual marginal, gives the result.
\end{proof}

\begin{theorem}[Logarithmic cost of a single crossing in a nested chain]
    \label{thm:chain_term_conjugation_logarithmic}
    \lean{Entropy.abs_re_inner_sub_chain_conj_le_logarithmic}
    \leanok
    Let $D_0\subseteq\cdots\subseteq D_{m-1}$ be regions of a finite
    tensor product, and let $\varphi$ be a unit vector. Suppose that
    each positive invertible filter $K_j$ on $D_j$ has a common
    orthonormal eigenbasis with $\rho_{\varphi,D_j}$, with eigenvalues
    $l_{j,i}>0$ and $p_{j,i}\ge0$. Let $a_j\ge0$. For each $j$,
    assume that for every pair $i,k$ with $p_{j,i},p_{j,k}>0$,
    \begin{align*}
        |\log l_{j,i}-\log l_{j,k}|
        \le\frac{a_j}{2}|\log p_{j,i}-\log p_{j,k}|.
    \end{align*}
    Let $X$ be Hermitian, supported on $S$, with
    $\lVert X\rVert\le c_0$, and put
    \begin{align*}
        d_S          & =\prod_{v\in S}n_v,\qquad \ell_S=\log(\mathrm e d_S), \\
        \mathcal J_S & =\{j:S\cap D_j\ne\varnothing,
        \ S\setminus D_j\ne\varnothing\}.
    \end{align*}
    Assume $|\mathcal J_S|\le1$ and
    $2a_j\ell_S\le1$ for every $j\in\mathcal J_S$.
    With $\widehat K_j=K_j\otimes I_{D_j^c}$ and
    $K=\widehat K_{m-1}\cdots\widehat K_0$, one has
    \begin{align*}
        \left|\operatorname{Re}\left(
        \langle\varphi,X\varphi\rangle
        -\langle\varphi,KXK^{-1}\varphi\rangle
        \right)\right|
        \le\sum_{j\in\mathcal J_S}
        8\mathrm e\,c_0\ell_S^2a_j^2.
    \end{align*}
    No smallness condition is imposed at indices outside
    $\mathcal J_S$.
\end{theorem}
\begin{proof}\leanok
    \uses{thm:split_filter_conjugation_logarithmic, lem:local_lift}
    If $j\in\mathcal J_S$, nesting and the single-crossing assumption
    imply that every earlier region is disjoint from $S$ and every
    later region contains $S$. The earlier filters commute with $X$.
    The later filters commute with their actual marginals, and the
    trace identity of Lemma~\ref{lem:local_lift} removes them from
    the expectation, starting with the outermost filter. The remaining
    expectation difference is the conjugation by $\widehat K_j$.
    Apply Theorem~\ref{thm:split_filter_conjugation_logarithmic}.
    If $\mathcal J_S$ is empty, the same disjoint-support and trace
    identities leave the expectation unchanged, and the sum is zero.
\end{proof}

\begin{theorem}[Sharp stationary-chain excitation energy]
    \label{thm:patch_sharp_energy}
    \lean{Entropy.re_inner_sub_le_of_chain_sharp}
    \leanok
    Under the hypotheses of Theorem~\ref{thm:patch_sharp_chain}, let
    $H=\sum_i h_i$, where each $h_i$ is Hermitian, supported on $S_i$,
    has norm at most $J$, and is split by at most one region. If
    $KHK^{-1}\phi=E_0\phi$, then
    \begin{align*}
        \operatorname{Re}\langle\phi,H\phi\rangle-E_0
        \le \sum_j a_j^2\sum_{i:\,D_j\text{ splits }S_i}
        d_{S_i\cap D_j}^2 J.
    \end{align*}
\end{theorem}
\begin{proof}\leanok
    \uses{thm:patch_sharp_chain}
    Expand the expectation difference as a sum over $i$, bound each
    summand by its absolute value, and apply
    Theorem~\ref{thm:patch_sharp_chain}. Interchange the finite sums to
    group the interactions by their crossing region. The eigenvector
    equation and unit normalization identify the conjugated expectation
    with $E_0$.
\end{proof}

These estimates assume the stated commutation and clipping properties.
Deriving them from a particular minimum and counting physical crossing
edges are separate steps in the patch argument.

\begin{theorem}[Logarithmic energy of a stationary nested product]
    \label{thm:chain_logarithmic_energy}
    \lean{Entropy.re_inner_sub_le_of_chain_logarithmic}
    \leanok
    Let $\mathcal H=\bigotimes_{v\in V}\mathbb C^{n_v}$ be a finite
    tensor product, let $D_0\subseteq\cdots\subseteq D_{m-1}$ be
    regions, and let $\varphi$ be a unit vector. Suppose that each
    positive invertible filter $K_j$ on $D_j$ has a common
    orthonormal eigenbasis with the actual marginal
    $\rho_{\varphi,D_j}$, with eigenvalues $l_{j,r}>0$ and
    $p_{j,r}\ge0$. Let $a_j\ge0$. For each $j$ and each pair $r,s$
    with $p_{j,r},p_{j,s}>0$, assume
    \begin{align*}
        |\log l_{j,r}-\log l_{j,s}|
        \le\frac{a_j}{2}|\log p_{j,r}-\log p_{j,s}|.
    \end{align*}
    Let $H=\sum_{i\in I}h_i$ be a finite sum of Hermitian operators,
    with $h_i$ supported on $S_i$ and $\lVert h_i\rVert\le c_0$.
    Put
    \begin{align*}
        d_i          & =\prod_{v\in S_i}n_v,\qquad \ell_i=\log(\mathrm e d_i), \\
        \mathcal C_j & =\{i\in I:S_i\cap D_j\ne\varnothing,
        \ S_i\setminus D_j\ne\varnothing\}.
    \end{align*}
    Assume that each $S_i$ is split by at most one region $D_j$.
    For every $j$ and every $i\in\mathcal C_j$, assume
    \begin{align*}
        2a_j\ell_i\le1.
    \end{align*}
    Write $\widehat K_j=K_j\otimes I_{D_j^c}$ and
    $K=\widehat K_{m-1}\cdots\widehat K_0$. If $E_0$ is real and
    $KHK^{-1}\varphi=E_0\varphi$, then
    \begin{align*}
        \operatorname{Re}\langle\varphi,H\varphi\rangle-E_0
        \le8\mathrm e\sum_{j=0}^{m-1}a_j^2
        \sum_{i\in\mathcal C_j}c_0\ell_i^2.
    \end{align*}
\end{theorem}
\begin{proof}\leanok
    \uses{thm:chain_term_conjugation_logarithmic}
    Unit normalization and the conjugated eigenvector equation give
    $\langle\varphi,KHK^{-1}\varphi\rangle=E_0$. Thus the energy
    difference is the sum, over $i$, of the real expectation
    differences for $h_i$. Bound each summand by its absolute value
    and apply Theorem~\ref{thm:chain_term_conjugation_logarithmic}.
    The smallness assumption applies precisely at the region that
    splits $S_i$, if such a region exists. Interchanging the finite
    sums gives the stated expression. The coefficient is that of
    \href{https://github.com/LionSR/TNLean/blob/67ea93e29f7c088f766b271208a2fde4cccf04a6/docs/notes/one_dimensional_note/part1_arealaw.tex#L711-L734}{Remark~6.3 of the one-dimensional note};
    the summation is the passage to the product energy in the
    area-law manuscript, \texttt{02-initial.tex}, lines 407--441.
\end{proof}

\begin{theorem}[Logarithmic energy of a maximizing nested product]
    \label{thm:filter_max_logarithmic_energy}
    \lean{Entropy.re_inner_sub_le_of_max_logarithmic}
    \leanok
    Let $D_0\subseteq\cdots\subseteq D_{m-1}$ be regions of a finite
    tensor product $\mathcal H=\bigotimes_{v\in V}\mathbb C^{n_v}$.
    Let $a_j>0$ and $f_j>0$, with
    $f_j\dim\mathcal H_{D_j}<1$ for every $j$. Suppose that each
    $K_j$ is feasible at floor $f_j$ and weight $a_j$, and that the
    family maximizes the norm of its output on a nonzero vector
    $\Omega$ among all such feasible families. Here feasibility means
    that $K_j=U_j\operatorname{diag}(y_{j,r}^{a_j/2})U_j^*$ for a
    unitary $U_j$ and real weights $y_{j,r}\ge f_j$ summing to one.
    Let $H=\sum_{i\in I}h_i$ be a finite sum of Hermitian operators,
    with $h_i$ supported on $S_i$ and $\lVert h_i\rVert\le c_0$.
    Assume that $H\Omega=E_0\Omega$ for a real $E_0$, and that each
    $S_i$ is split by at most one $D_j$. Define
    \begin{align*}
        d_i          & =\prod_{v\in S_i}n_v,\qquad \ell_i=\log(\mathrm e d_i), \\
        \mathcal C_j & =\{i\in I:S_i\cap D_j\ne\varnothing,
        \ S_i\setminus D_j\ne\varnothing\}.
    \end{align*}
    For every $j$ and every $i\in\mathcal C_j$, assume
    \begin{align*}
        2a_j\ell_i\le1.
    \end{align*}
    Write $\widehat K_j=K_j\otimes I_{D_j^c}$ and put
    \begin{align*}
        K & =\widehat K_{m-1}\cdots\widehat K_0,\qquad
        \psi=K\Omega,\qquad \varphi=\frac{\psi}{\lVert\psi\rVert}.
    \end{align*}
    Then $\psi\ne0$, and the actual normalized output satisfies
    \begin{align*}
        \operatorname{Re}\langle\varphi,H\varphi\rangle-E_0
        \le8\mathrm e\sum_{j=0}^{m-1}a_j^2
        \sum_{i\in\mathcal C_j}c_0\ell_i^2.
    \end{align*}
\end{theorem}
\begin{proof}\leanok
    \uses{prop:filter_maximizer, thm:all_filters_clipped,
        thm:chain_logarithmic_energy}
    Positive floors make every $K_j$ invertible, so
    Proposition~\ref{prop:filter_maximizer} gives $\psi\ne0$ and
    $\varphi$ is a unit vector. Actual maximality gives marginal
    commutation and common-basis clipping for $\psi$ by
    Theorem~\ref{thm:all_filters_clipped}. Normalization multiplies
    each marginal by the same positive scalar
    $\lVert\psi\rVert^{-2}$. Commutation is preserved, and the
    added constant in each logarithmic marginal eigenvalue cancels
    in differences on the positive spectrum. The filters therefore
    satisfy the same clipping inequalities relative to the actual
    marginals of $\varphi$. The eigenvector equation gives
    $KHK^{-1}\varphi=E_0\varphi$. Apply
    Theorem~\ref{thm:chain_logarithmic_energy}. These are the
    optimizer and normalization passages in the area-law manuscript,
    \texttt{02-initial.tex}, lines 342--441; no clipping or
    normalization data of an unrelated vector are assumed.
\end{proof}

\begin{theorem}[Sharp energy of a maximizing nested product]
    \label{thm:filter_max_sharp_energy}
    \lean{Entropy.re_inner_sub_le_of_max_sharp,
        Entropy.abs_re_inner_sub_le_of_max_sharp}
    \leanok
    \uses{prop:stationary_filter_clipped,lem:local_lift}
    Let $D_0\subseteq\cdots\subseteq D_{m-1}$ be regions of a finite
    system, let $0<a_j\le1$, and let $f_j>0$ with
    $f_j\dim\mathcal H_{D_j}<1$. Suppose a feasible family $K_j$
    maximizes the output norm. Assume $\Omega\ne0$,
    $H\Omega=E_0\Omega$, and $H=\sum_i h_i$, where each $h_i$ is
    Hermitian, supported on $S_i$, and has norm at most $c_0$.
    Each $S_i$ is split by at most one $D_j$. Put
    \begin{align}
        B_{\mathrm{in}}
         & =\sum_j a_j^2\sum_{i:\,D_j\text{ splits }S_i}
        (\dim\mathcal H_{S_i\cap D_j})^2c_0.
        \label{eq:filter_inside_budget}
    \end{align}
    The normalized output $\phi=K_{m-1}\cdots K_0\Omega/
        \lVert K_{m-1}\cdots K_0\Omega\rVert$ satisfies
    $|\operatorname{Re}\langle\phi,H\phi\rangle-E_0|
        \le B_{\mathrm{in}}$, and hence also the corresponding one-sided bound.
\end{theorem}
\begin{proof}\leanok
    \uses{thm:all_filters_clipped,prop:stationary_filter_clipped,
        prop:filter_maximizer,thm:chain_energy,thm:patch_sharp_energy}
    Positive floors make the output nonzero. Clipping and commutation
    hold for every filter by Theorem~\ref{thm:all_filters_clipped} and
    persist after normalization. The conjugated eigenvector equation is
    $KHK^{-1}\phi=E_0\phi$, so Theorem~\ref{thm:patch_sharp_energy}
    gives the upper bound with (\ref{eq:filter_inside_budget}).
    Apply the same argument to $-h_i$ and eigenvalue $-E_0$.
    The feasible family and its maximizing property are unchanged;
    supports, Hermiticity and norm bounds are preserved. The second
    inequality bounds the negative of the same energy difference.
\end{proof}

\begin{theorem}[Sharp absolute energy without a positive floor]
    \label{thm:filter_zero_floor_sharp_energy}
    \lean{Entropy.exists_max_feasible_chain_zero_floor_abs_energy_sharp,
        Entropy.exists_max_feasible_chain_zero_floor_normalized_abs_energy_sharp}
    \leanok
    \uses{prop:stationary_filter_clipped,lem:local_lift}
    Let $D_0\subseteq\cdots\subseteq D_{m-1}$ and $0<a_j\le1$.
    Suppose $\Omega\ne0$, $H\Omega=E_0\Omega$, and $H=\sum_i h_i$,
    with $h_i$ Hermitian, supported on $S_i$, and
    $\lVert h_i\rVert\le c_0$. Assume each $S_i$ is split by at
    most one $D_j$, and let $B_{\mathrm{in}}$ be
    (\ref{eq:filter_inside_budget}). There is a floor-zero feasible
    family $K_j$ that maximizes the output norm and whose output
    $\psi=K_{m-1}\cdots K_0\Omega$ is nonzero and satisfies
    \begin{align}
        |\operatorname{Re}\langle\psi,H\psi\rangle-E_0\lVert\psi\rVert^2|
         & \le B_{\mathrm{in}}\lVert\psi\rVert^2.
        \label{eq:filter_zero_sharp_energy}
    \end{align}
    The same family has unit normalized output
    $\phi=\psi/\lVert\psi\rVert$ and
    $|\operatorname{Re}\langle\phi,H\phi\rangle-E_0|
        \le B_{\mathrm{in}}$. The filters may be singular.
\end{theorem}
\begin{proof}\leanok
    \uses{thm:filter_max_sharp_energy,thm:filter_zero_floor_closed_max,
        lem:filter_zero_floor_nonzero,prop:filter_maximizer}
    At each positive floor, Theorem~\ref{thm:filter_max_sharp_energy}
    gives (\ref{eq:filter_zero_sharp_energy}) after multiplication
    by the squared output norm. The absolute homogeneous difference
    is continuous in the filter matrices, including at zero output.
    Its sublevel condition is therefore closed. Apply
    Theorem~\ref{thm:filter_zero_floor_closed_max} to obtain a
    maximizing family satisfying that condition.
    Lemma~\ref{lem:filter_zero_floor_nonzero} proves nonvanishing,
    after which division by the squared norm gives the normalized bound.
\end{proof}
